
\documentclass[10pt,twocolumn,letterpaper]{article}

\usepackage{conf}
\usepackage{times}
\usepackage{inconsolata}
\usepackage[utf8]{inputenc}
\usepackage[T1]{fontenc}

\usepackage{microtype}
\usepackage{graphicx}
\usepackage{subcaption}
\usepackage{booktabs}
\usepackage{tabularx}
\newcolumntype{Y}{>{\centering\arraybackslash}X}

\usepackage[round]{natbib}

\usepackage{algorithm}

\usepackage[table,dvipsnames]{xcolor}
\usepackage[pagebackref=true,breaklinks=true,colorlinks,bookmarks=false]{hyperref}
\definecolor{mydarkblue}{rgb}{0,0.08,0.45}
\hypersetup{
  linkcolor=mydarkblue,
  citecolor=mydarkblue,
  filecolor=mydarkblue,
  urlcolor=mydarkblue,
}

\usepackage{amsmath}
\usepackage{amssymb}
\usepackage{mathtools}
\usepackage{amsthm}

\usepackage[capitalize,noabbrev]{cleveref}

\theoremstyle{plain}
\newtheorem{theorem}{Theorem}[section]
\newtheorem{proposition}[theorem]{Proposition}
\newtheorem{lemma}[theorem]{Lemma}
\newtheorem{corollary}[theorem]{Corollary}
\theoremstyle{definition}
\newtheorem{definition}[theorem]{Definition}
\newtheorem{assumption}[theorem]{Assumption}
\theoremstyle{remark}
\newtheorem{remark}[theorem]{Remark}

\usepackage[disable,textsize=tiny]{todonotes}


\usepackage{algpseudocode}
\usepackage{xspace}

\newcommand{\flashOptim}{FlashOptim\xspace}
\newcommand{\flashAdam}{FlashAdamW\xspace}
\newcommand{\flashSGD}{FlashSGD\xspace}
\newcommand{\flashLion}{FlashLion\xspace}

\definecolor{ultramarine}{RGB}{0,32,96}

\newcommand{\hlcolor}{SpringGreen}  
\newcommand{\hlchange}[1]{{{\fboxsep=1pt\text{\colorbox{\hlcolor}{$#1$}}}}}
\newcommand{\hltext}[1]{\colorbox{\hlcolor}{#1}}


\definecolor{firebrick}{RGB}{178,34,34}
\definecolor{navy}{RGB}{0,0,128}
\definecolor{forestgreen}{RGB}{0,128,0}

\ifdefined\final
  \newcommand{\authorcomment}[3]{}
  \newcommand{\draft}[1]{}
  \newcommand{\myinclude}[1]{\input{#1}}
\else
  \newcommand{\authorcomment}[3]{{\color{#2}[\textbf{#1}:#3]}}
  \newcommand{\draft}[1]{{\color{ultramarine}\itshape #1}}
  \newcommand{\myinclude}[1]{\input{#1}}
\fi

\newcommand{\JJ}[1]{\authorcomment{JJ}{firebrick}{#1}}
\newcommand{\DB}[1]{\authorcomment{DB}{forestgreen}{#1}}
\newcommand{\AG}[1]{\authorcomment{AG}{navy}{#1}}

\usepackage{mathtools}
\DeclarePairedDelimiter{\abs}{\lvert}{\rvert}

\usepackage{amsmath}
\DeclareMathOperator{\signop}{sign}
\newcommand{\sign}[1]{\signop\!\left(#1\right)}

\newcommand{\thetalp}{\theta^\prime}


\algtext*{EndFunction}
\algtext*{EndProcedure}
\algtext*{EndIf}
\algtext*{EndFor}
\algtext*{EndWhile}
\algtext*{EndLoop}
\algtext*{EndRepeat}

\providecommand*{\algorithmautorefname}{Algorithm}
\providecommand*{\subfigureautorefname}{Figure}  %
\algrenewcommand{\algorithmiccomment}[1]{// #1}



\newcommand{\nseed}{3\xspace}

\usepackage{enumitem}

\usepackage{siunitx}
\sisetup{exponent-mode=scientific, round-mode=figures, round-precision=3}
\sisetup{exponent-product=\cdot}

\usepackage{setspace}



\conffinalcopy

\begin{document}

\title{FlashOptim: Optimizers for Memory Efficient Training}

\author{Jose Javier Gonzalez Ortiz\thanks{Equal contribution}\\
Databricks AI Research\\
{\tt\small j.gonzalez@databricks.com}
\and
Abhay Gupta\\
Databricks AI Research\\
{\tt\small abhay.gupta@databricks.com}
\and
Chris Renard\\
Databricks AI Research\thanks{Now at Standard Kernel Co.}\\
{\tt\small chris@standardkernel.co}
\and
Davis Blalock\footnotemark[1]\\
Databricks AI Research\thanks{Now at Google DeepMind}\\
{\tt\small daviswblalock@gmail.com}
}

\maketitle

\begin{abstract}

Standard mixed-precision training of neural networks requires many bytes of accelerator memory for each model parameter. These bytes reflect not just the parameter itself, but also its gradient and one or more optimizer state variables. With each of these values typically requiring 4 bytes, training even a 7 billion parameter model can be impractical for researchers with less than 100GB of accelerator memory.

We introduce \flashOptim, a suite of optimizations that reduces per-parameter memory by over 50\% while preserving model quality and API compatibility.
Our approach introduces two key techniques. First, we improve master weight splitting by finding and exploiting a tight bound on its quantization error. Second, we design companding functions that greatly reduce the error in 8-bit optimizer state quantization.
Together with 16-bit gradients, these techniques reduce AdamW memory from 16 bytes to 7 bytes per parameter, or 5 bytes with gradient release. They also cut model checkpoint sizes by more than half.

Experiments with \flashOptim applied to SGD, AdamW, and Lion show no measurable quality degradation on any task from a collection of standard vision and language benchmarks, including Llama-3.1-8B finetuning.

\end{abstract}


\section{Introduction} \label{sec:intro}

Recent advances in deep learning have been driven largely by scaling: larger models trained on more data consistently yield better results across language~\citep{kaplan2020scaling,hoffmann2022training,chowdhery2023palm} and vision~\citep{rosenfeld2019constructive,tan2019efficientnet,dehghani2023scalingvit} domains.
Training large models can require a great deal of accelerator memory, with each training iteration requiring memory to store parameters, activations, gradients, and optimizer state.

How much memory do these tensors require? \autoref{tbl:bytesPerParam} shows a typical breakdown. Excluding activations, which scale with batch size rather than parameter count, training with Adam uses about 16 bytes per parameter. Training a 7-billion-parameter LLM therefore requires at least 112GB of accelerator memory, plus more memory for activations.

\begin{table}[t]
    \centering
    \caption{\textbf{Memory per parameter (bytes) for model training}. 
    \flashOptim reduces Adam from 16 to 7 bytes and SGD from 12 to 6 bytes.
     $(\star)$ With gradient release, we further reduce total memory requirements by 2 bytes.}
    \label{tbl:bytesPerParam}
    \setlength{\tabcolsep}{3pt}

    \small
    \rowcolors{2}{white}{gray!10}
    \renewcommand{\arraystretch}{1.2}
    \begin{tabular}{lcccc}
        \toprule
        \textbf{Tensor} & \textbf{SGD} & \textbf{\flashSGD} & \textbf{Adam} & \textbf{FlashAdam} \\
        \midrule
        Master Weights      & 4 & 2\phantom{ (0${}^\star$)} & 4 & 2\phantom{ (0${}^\star$)} \\
        Weight Correction   &   & 1\phantom{ (0${}^\star$)} &   & 1\phantom{ (0${}^\star$)} \\
        Gradients           & 4 & 2 (0${}^\star$) & 4 & 2 (0${}^\star$) \\
        Momentum            & 4 & 1\phantom{ (0${}^\star$)} & 4 & 1\phantom{ (0${}^\star$)} \\
        Variance            &   &   & 4 & 1\phantom{ (0${}^\star$)} \\
        \midrule
        \textbf{Total} & \textbf{12} & \textbf{6 (4${}^\star$)} & \textbf{16} & \textbf{7 (5${}^\star$)} \\
        \bottomrule
    \end{tabular}
\end{table}

Several approaches mitigate this memory consumption. Distributed training with tensor sharding~\citep{rajbhandari2020zero} divides the memory load across multiple accelerators. While this is standard practice in well-resourced organizations, it requires access to multiple accelerators that many practitioners lack.
Another alternative is to perform CPU offloading~\citep{ren2021zero}, which moves some tensors to host memory at the cost of added overhead and complexity. Third, parameter-efficient methods~\citep{li2021prefix,hu2022lora} reduce trainable parameters by freezing most weights and training either a small subset of the original weights or a small set of new auxiliary weights, but fundamentally alter the training dynamics~\citep{biderman2024lora}.


In this work, we describe \flashOptim, a set of techniques to reduce 
parameter-associated memory in common deep learning optimizers. 
\autoref{fig:finetuning_memory_breakdown} shows an example: with \flashOptim, finetuning 
Llama-3.1-8B drops from 175 GiB to 113 GiB peak memory.
Crucially, these memory savings are effectively free---\flashOptim runs just as fast as standard optimizers and causes no measurable loss of model quality across a suite of established training tasks (\S\ref{sec:results}). This allows our optimizer implementations to serve as drop-in replacements for their unoptimized counterparts. \flashOptim incorporates existing enhancements, such as gradient release~\citep{zhang2023adam,warner2024optimi}, while also introducing improved float splitting~\citep{zamirai2020revisiting,warner2024optimi} and simplified 8-bit optimizer state quantization~\citep{dettmers2022optimizers,peng2023fp8lm,xi2025coat,fishman2024scaling}. Furthermore, \flashOptim composes cleanly with existing memory-reduction techniques, such as sharding tensors across accelerators, offloading to CPU, or freezing parameters.

\begin{figure}
    \centering
    \includegraphics[width=\columnwidth]{figures/finetuning_memory_breakdown.pdf}
    \caption{\textbf{Memory breakdown for finetuning Llama-3.1-8B.} \flashOptim reduces peak memory from 175 to 113 GiB by compressing parameters and optimizer states.}
    \label{fig:finetuning_memory_breakdown}
\end{figure}

We make the following contributions:

\begin{itemize}[itemsep=-2pt]
    \item \textbf{Improved float splitting}: Instead of materializing both a 32-bit master weight and a 16-bit downcast weight for forward and backward, one can split each master weight into a low-precision weight and a correction term stored in the optimizer~\citep{zamirai2020revisiting,warner2024optimi}. We improve on existing float splitting techniques by (a) enabling either 8- or 16-bit error correction and (b) achieving much lower reconstruction error for a given number of correction bits. This allows us to use 24-bit master weights with no loss of model quality.
    \item \textbf{Companded optimizer state quantization}: Several works have shown that one can compress optimizer states to 8 bits per element given sufficient software complexity. We demonstrate that one can do this much more simply, with nothing more than a one-line preprocessing function before standard group-wise linear quantization. Our ablations across different tensor types suggest that designing custom companding functions is a fruitful direction for future research. 
    \item \textbf{Fused optimized kernels}: We implement \flashOptim as optimizer step kernels that fuse all compression and quantization operations, reducing memory while preserving throughput during training. Our implementation is publicly available at \url{https://github.com/databricks/flashoptim}.
\end{itemize}






\section{Related Work} \label{sec:relatedWork}

\textbf{Low-Precision Training.}
Mixed-precision training~\citep{micikevicius2018mixed} executes forward and backward passes in FP16 to reduce memory and compute, while retaining FP32 precision for optimizer states and master weights to preserve numerical stability.
\citet{kalamkar2019bfloat16} showed that BFloat16~\citep{google_cloud_bfloat16_2019} works equally well, and \citet{zamirai2020revisiting} explored pure BF16 master weights with stochastic rounding and Kahan summation.
Recent work has pushed further with FP8 training~\citep{wang2018training,mellempudi2019mixed,micikevicius2022fp8,fishman2024scaling,narayan2025mu}, though these approaches primarily target compute formats and retain higher-precision storage for master weights.
\flashOptim extends this line of work with an improved float splitting mechanism that reduces storage to 3 bytes per parameter, down from 4-byte FP32, while maintaining FP32-equivalent training semantics.


\textbf{Optimizer State Compression.}
\citet{dettmers2022optimizers} applied 8-bit block-wise dynamic quantization to Adam's momentum and variance, reducing optimizer state from 8 to 2 bytes per parameter.
Follow-up work explored FP8 representations~\citep{peng2023fp8lm,xi2025coat,fishman2024scaling}, and \citet{li2023memory} compressed both moments to 4-bit using row and column-wise quantization.
MicroAdam~\citep{modoranu2024microadam} instead compresses gradients before updating optimizer states. Rather than design elaborate quantization methods or number formats, we show that one can obtain 8-bit optimizer states with no quality loss using simple, one-line preprocessing functions.
Beyond optimizer states, we address additional sources of per-parameter memory, eliminating entire bytes from other tensors.


\textbf{Gradient Memory and Communication.}
LOMO~\citep{lv2024lomo}, AdaLOMO~\citep{lv2024adalomo}, and Adam Accumulation~\citep{zhang2023adam} fuse parameter updates into the backward pass to release gradient memory eagerly. However, this conflicts with gradient accumulation, which requires the full accumulated gradient before updating.
In distributed settings, gradient communication can also become a bottleneck. One can reduce this bottleneck by, e.g., compressing gradients to 1-bit with error feedback~\citep{tang20211bit}, or using low-rank approximations~\citep{vogels2019powersgd}.
\flashOptim supports gradient release when compatible and could be used alongside communication compression techniques.


\textbf{Memory-Efficient Optimization.}
Alternative optimizer designs reduce memory by restructuring update rules and stored buffers.
Adafactor~\citep{shazeer2018adafactor} achieves sublinear memory by factorizing the second moment into row and column statistics; SM3~\citep{anil2019memory} stores structured maxima; NovoGrad~\citep{ginsburg2019stochastic} replaces per-parameter variance with layer-wise normalization.
Adam-mini~\citep{zhang2025adammini} shares variance terms across parameter blocks, while Adapprox~\citep{zhao2024adapprox} uses a low-rank approximation.
Other approaches eliminate the second moment entirely: Lion~\citep{chen2023lion} uses sign-based momentum, and Muon~\citep{jordan2024muon,liu2025muon} applies orthogonalized updates.
\citet{pethick2025training} extend Muon to unify gradient accumulation with momentum, removing dedicated optimizer memory altogether.

Low-rank decompositions approximate full tensors while requiring less memory. 
For fine-tuning, LoRA~\citep{hu2022lora} and QLoRA~\citep{dettmers2023qlora} freeze base weights and train only low-rank adapters.
For pretraining, GaLore~\citep{zhao2024galore} projects gradients to a low-rank subspace, and APOLLO~\citep{zhu2025apollo} approximates adaptive scaling with random projections.
Unlike these approaches that modify the optimizer's update rule, \flashOptim preserves standard optimizer semantics and can be combined with these techniques.


\textbf{System-Level Memory Optimizations.}
System-level approaches reduce accelerator memory without changing optimization semantics.
Activation checkpointing~\citep{chen2016training,korthikanti2023reducing} trades compute for memory by recomputing activations during the backward pass.
ZeRO~\citep{rajbhandari2020zero} partitions optimizer states, gradients, and parameters across data-parallel ranks, while offloading~\citep{rajbhandari2021zeroinfinity,ren2021zero} moves state to CPU or NVMe memory.
\flashOptim is orthogonal to these approaches: it reduces the per-rank footprint and can be used with ZeRO, FSDP~\citep{zhao2023pytorchfsdp}, and activation checkpointing.



\section{Method} \label{sec:method}

This section describes the two key techniques behind \flashOptim: weight splitting (\S\ref{sub:ulp_weight_compression}) and companded optimizer state quantization (\S\ref{sub:optimizer_state_quantization_with_companding}).
We then describe how to integrate these ideas into common optimizer updates while minimizing associated overhead (\S\ref{sub:optimizer_update}).



\subsection{Weight Splitting}
\label{sub:ulp_weight_compression}

Mixed-precision training uses 16-bit weights for forward and backward passes, but accumulating gradient updates requires higher precision to avoid stagnation~\citep{micikevicius2018mixed}.
Thus, the standard practice is to maintain FP32 precision master weights during training.

However, this introduces waste: the downcast weights take space, but store no information beyond what's saved in the master weights. To eliminate this redundancy, weight splitting~\citep{zamirai2020revisiting,warner2024optimi} instead stores the downcast weights and narrow error-correction values. By combining a 16-bit weight $\thetalp$ with a 16-bit error correction value $\rho$, one has enough information to reconstruct a 32-bit master weight $\theta$ with no redundancy. 


The core questions in a weight splitting scheme are 1) how to set $(\thetalp, \rho)$ given $\theta$ and 2) how to estimate $\theta$ given $(\thetalp, \rho)$.  
One obvious approach is to use the high 16 bits of $\theta$ as $\thetalp$ and the low 16 bits as $\rho$. This admits exact reconstruction of $\theta$ for the special case of BF16 $\thetalp$ and FP32 $\theta$, since these formats happen to share the same exponent sizes and offsets. However, this approach doesn't generalize to other pairs of number formats. It also rounds towards zero instead of towards the nearest low-precision value. 

A more general alternative, used in previous work \citep{zamirai2020revisiting,warner2024optimi}, is to set $\rho = \theta - \thetalp$, represented as a BF16 value. The problem with this is that the difference of two floating-point numbers requires as many bits to store exactly as the wider of the two floats.\footnote{Consider, e.g., storing the minimum float32 subnormal, which will be rounded to zero by any narrower datatype.} This means that a BF16 $\rho$ incurs approximation error. E.g., if the rounding error were 1e-5, BF16 could only represent 1.0014e-5. In general, while BF16's wide exponent lets it represent nearly the full range of FP32, its 7 mantissa bits only guarantee a relative error bound of 1/256.  

Our observation is that \textbf{all exponent bits in this scheme are wasted}. 
The exponent of $e \triangleq \theta - \thetalp$ can always be inferred from 
$\thetalp$: under round-to-nearest, $\theta$ must lie within 
$[\thetalp - u/2, \thetalp + u/2]$, where $u = \text{ULP}(\thetalp)$ is the 
unit in the last place~\citep{goldberg1991every}. If $\theta$ were outside 
this interval, it would have rounded to a different value. It therefore 
suffices to encode where $e$ falls within this tiny interval rather than 
across the full FP32 range.

To exploit this observation, we rescale $e$ such that $[-u/2, u/2]$ maps to $[-N, N]$; $N \triangleq 2^b -1$ and then quantize this rescaled $e$ to the nearest $b$-bit integer. That is,
\begin{equation}
\begin{aligned}
    \thetalp &= \text{downcast}(\theta) \\
    \rho &= \text{round}\left( \frac{\theta - \thetalp}{\text{ULP}(\thetalp)/2}\cdot N \right), \\
\end{aligned}
\label{eq:compress}
\end{equation}

To reconstruct $\theta$ from $(\thetalp, \rho)$, we invert this scaling and add the result to $\thetalp$. 
\begin{equation}
\begin{aligned}
\hat{\theta} =    \thetalp + \frac{\rho}{N} \cdot \frac{\text{ULP}(\thetalp)}{2}
\end{aligned}
\label{eq:decompress}
\end{equation}
For tensors of floating-point values we apply this transformation elementwise.
\autoref{alg:split_reconstruct} provides a lower-level description of the compression and decompression operations with numerical precision considerations.

\begin{algorithm}[h]
    \caption{Weight Splitting}
    \label{alg:split_reconstruct}
    \begin{algorithmic}[1]
    \Statex \textbf{Constants:} $N$ (127 for INT8, 32767 for int16)
    \Function{$\mathcal{C}$}{$\theta$}
        \State $\thetalp \gets \mathrm{Downcast}(\theta)$
        \State $e \gets \theta - \mathrm{Float32}(\thetalp)$
        \State $\ell \gets \lfloor\mathrm{log_2ULP}(\thetalp)\rfloor - 1$
        \State $h \gets \lfloor -\ell / 2 \rfloor$ \Comment{For numerical stability}
        \State $e_{\mathrm{norm}} \gets (e \cdot 2^h) \cdot 2^{-\ell - h}$
        \State $\rho \gets \mathrm{Int}(\mathrm{Round}(\mathrm{Clamp}(e_{\mathrm{norm}}, -1, 1) \cdot N))$
        \State \Return $\thetalp, \rho$
    \EndFunction
    \vspace{0.25\baselineskip}
    \Function{$\mathcal{C}^{-1}$}{$\thetalp, \rho$} 
        \State $\ell \gets \lfloor\mathrm{log_2ULP}(\thetalp)\rfloor - 1$
        \State $h \gets \lfloor \ell / 2 \rfloor$
        \State $e \gets ((\mathrm{Float32}(\rho) / N) \cdot 2^h) \cdot 2^{\ell - h}$
        \State \Return $\mathrm{Float32}(\thetalp) + e$
    \EndFunction
    \end{algorithmic}
    \end{algorithm}
    


When using BF16 for $\thetalp$ and INT8 for $\rho$, the compressed representation provides approximately 24 bits of effective precision (16 from BF16 plus 8 from the error term).
This is analogous to the PXR24 format used in high-dynamic-range imaging, which achieves similar precision by rounding 32-bit floats to 24 bits~\citep{kainz2004openexr}.

\subsection{Companded Optimizer State Quantization}
\label{sub:optimizer_state_quantization_with_companding}

For optimizer state variables such as momentum and variance estimates, a common approach is group-wise quantization: dividing tensors into fixed-length groups and mapping values to a lower-precision format like INT8~\citep{dettmers2022optimizers}.
To increase the precision range of the group of values, they are rescaled using the maximum absolute value (\emph{absmax}), which is stored as an additional scale with 32 or 16 bits of precision.
While simple, this \emph{uniform} quantization allocates bins evenly across the value range, implicitly assuming that values are roughly uniformly distributed.
Our measurements show that optimizer state distributions violate this assumption, and we find that applying nonlinear \emph{companding} functions before quantization can reshape these distributions toward uniformity, improving utilization of quantization bins and reducing quantization error.
As we show in \S\ref{sec:opt_quant}, this companding step is critical: without it, linear quantization of optimizer states causes training to diverge.

We design specialized transformations for each optimizer state type.
For momentum tensors (used in SGD, Adam, and Lion), we first normalize each group by its absmax scale, then apply a softsign-like function:
\begin{equation}
    \phi_m(x) = \frac{2x}{1+|x|} \qquad \phi_m^{-1}(z) = \frac{z}{2-|z|}
\end{equation}
This function compresses extreme values: inputs near $\pm 1$ are pushed toward the center, spreading the momentum distribution more evenly across quantization bins.
In contrast, for variance tensors in Adam, we first apply a square root, then normalize by the group absmax:
\begin{equation}
    \phi_v(x) = \sqrt{x} \qquad \phi_v^{-1}(z) = z^2
\end{equation}
Here the square root is motivated by Adam's variance update
$v_t = \beta_2 v_{t-1} + (1 - \beta_2) g^2$ that accumulates squared gradients, producing heavy-tailed distributions with large dynamic range.
Both transformations satisfy key design criteria: they are exactly invertible, computationally efficient (one division or square root per element), their inverses are computationally efficient, and require no hyperparameters.

    \begin{algorithm}[t]
        \caption{$\mathcal{Q}_m$: Momentum Quantization}
        \label{alg:quant_mom}
        \begin{algorithmic}[1]
        \Statex \textbf{Constants:} Group size $G = 32$
        \vspace{0.25\baselineskip}
        \Function{$\mathcal{Q}_m$}{$m$}
            \For{each group $g$ of $G$ elements}
                \State $s_g \gets \max(|m_g|)$
                \State $m'_g \gets m_g / s_g$
                \State $m''_g \gets 2 m'_g / (1 + |m'_g|)$
                \State $m^q_g \gets \mathrm{Round}(m''_g \cdot 127, \text{INT8})$
            \EndFor
            \State \Return $m^q, s$
        \EndFunction
        \vspace{0.25\baselineskip}
        \Function{$\mathcal{Q}_m^{-1}$}{$m^q, s$}
            \For{each group $g$}
                \State $m''_g \gets m^q_g / 127$
                \State $m'_g \gets m''_g / (2 - |m''_g|)$
                \State $m_g \gets m'_g \cdot s_g$
            \EndFor
            \State \Return $m$
        \EndFunction
        \end{algorithmic}
        \end{algorithm}


\begin{algorithm}[t]
\caption{$\mathcal{Q}_v$: Variance Quantization}
\label{alg:quant_var}
\begin{algorithmic}[1]
\Statex \textbf{Constants:} Group size $G = 32$
\vspace{0.25\baselineskip}
\Function{$\mathcal{Q}_v$}{$v$}
    \State $v' \gets \sqrt{v}$
    \For{each group $g$ of $G$ elements}
        \State $s_g \gets \max(v'_g)$
        \State $v^q_g \gets \mathrm{Round}((v'_g / s_g) \cdot 255, \text{UINT8})$
    \EndFor
    \State \Return $v^q, s$
\EndFunction
\vspace{0.25\baselineskip}
\Function{$\mathcal{Q}_v^{-1}$}{$v^q, s$}
    \For{each group $g$}
        \State $v'_g \gets (v^q_g / 255) \cdot s_g$
        \State $v_g \gets (v'_g)^2$
    \EndFor
    \State \Return $v$
\EndFunction
\end{algorithmic}
\end{algorithm}


For both momentum and variance, we partition the tensor into groups of $G=32$ elements and store a separate FP16 scale factor per group, introducing an overhead of $2/G = 1/16$ bytes per parameter.
We store the normalized momentum in signed integers (INT8) and variance in unsigned integers (UINT8) since it is non-negative.
Algorithms~\ref{alg:quant_mom} and~\ref{alg:quant_var} detail the quantization and dequantization procedures for momentum and variance respectively.




\begin{algorithm}[t]
    \caption{FlashAdamW: Memory-Efficient AdamW. We \hlchange{\text{highlight}} the changes from AdamW.}
    \label{alg:flashadamw}
    \begin{algorithmic}[1]
    \Require Parameters $\theta_0$, learning rate schedule $\{\eta_t\}_{t=1}^{T}$, $\beta_1, \beta_2 \in [0,1)$, $\varepsilon > 0$, weight decay $\lambda \geq 0$, loss $\mathcal{L}(\theta)$, minibatch sampler $\mathcal{B}(\cdot)$
    \State $\hlchange{m_0^q, m_0^s \gets \mathcal{Q}_m(0)}$
    \State $\hlchange{v_0^q, v_0^s \gets \mathcal{Q}_v(0)}$
    \State $\hlchange{\thetalp_0, \rho_0 \gets \mathcal{C}(\theta_0)}$
    \State Initialize $t \gets 0$
    \While{not converged}
        \State $t \gets t + 1$
        \State $B_t \sim \mathcal{B}$
        \State $g_t \gets \nabla_\theta \mathcal{L}(B_t; \hlchange{\thetalp_{t-1}})$
        \State \Comment{Reconstruct optimizer state and master weight}
        \State $\hlchange{m_{t-1} \gets \mathcal{Q}_m^{-1}(m^q_{t-1}, m^s_{t-1})}$
        \State $\hlchange{v_{t-1} \gets \mathcal{Q}_v^{-1}(v^q_{t-1}, v^s_{t-1})}$
        \State $\hlchange{\theta_{t-1} \gets \mathcal{C}^{-1}(\thetalp_{t-1}, \rho_{t-1})}$
        \State \Comment{Standard optimizer update}
        \State $m_t \gets \beta_1 m_{t-1} + (1 - \beta_1) g_t$
        \State $v_t \gets \beta_2 v_{t-1} + (1 - \beta_2) g_t^2$
        
        \State $\hat{m}_t \gets m_t / (1 - \beta_1^t)$
        \State $\hat{v}_t \gets v_t / (1 - \beta_2^t)$
        
        \State $\theta_t \gets \theta_{t-1} - \eta_t \left( \hat{m}_t / (\sqrt{\hat{v}_t} + \varepsilon) + \lambda \theta_{t-1} \right)$
        \State \Comment{Quantize optimizer state and split master weight}
        \State $\hlchange{m^q_t, m^s_t \gets \mathcal{Q}_m(m_t)}$
        \State $\hlchange{v^q_t, v^s_t \gets \mathcal{Q}_v(v_t)}$
        \State $\hlchange{\thetalp_t, \rho_t \gets \mathcal{C}(\theta_t)}$
    \EndWhile
    \end{algorithmic}
    \end{algorithm}


\subsection{Optimizer Update}
\label{sub:optimizer_update}

We modify any given gradient update rule by adding a prologue and an epilogue. In the prologue, we dequantize the optimizer states and reconstruct the master weight $\theta$ from the low-precision weight $\thetalp$ and the error correction bits $\rho$. In the epilogue, we quantize the new optimizer state and split the new $\theta$ into an updated $(\thetalp, \rho)$. At the start of training, 
we downcast the master weights to BF16
to ensure training runs directly on the low-precision $\thetalp$ with no downcasts apart from our optimizer step.
\autoref{alg:flashadamw} illustrates these changes for the AdamW optimizer, and Algorithms~\ref{alg:flashsgd} and~\ref{alg:flashlion} in the appendix show the corresponding changes to SGD and Lion respectively.

\subsection{Implementation}

\textbf{Update kernels.}
Since compression and quantization are bandwidth-bound operations, we implement the optimizer step as a single fused Triton kernel~\citep{tillet2019triton}. 
For example, for the AdamW update our kernel encompasses steps 9 through 22 from~\autoref{alg:flashadamw}.

\textbf{Gradient release.}
We implement gradient release~\citep{zhang2023adam}, interleaving gradient computation with optimizer updates during backpropagation.
As each gradient is computed, we eagerly apply the optimizer rule to free the gradient memory.
We apply this optimization only when gradient accumulation is disabled. 

\textbf{Distributed training.}
Our implementation is compatible with parameter sharding approaches such as PyTorch FSDP~\citep{zhao2023pytorchfsdp}.
During forward and backward passes, only the 16-bit $\thetalp$ parameters are all-gathered; the correction term $\rho$ remains local with the optimizer states.

\textbf{Checkpoint size.}
Our representation reduces checkpoint size.
For instance, standard Adam checkpoints require 12 bytes per parameter (4 for weights, 4 for momentum, 4 for variance); \flashAdam requires only 5 bytes (2 for weights, 1 for correction, 1 for momentum, and 1 for variance).
For a 7B model, checkpoint size reduces from 84GB to 35GB.

\textbf{Code availability.}
We make our implementation widely available as an open-source PyTorch library at \url{https://github.com/databricks/flashoptim}. 






\section{Experiments} \label{sec:results}

\subsection{Experimental Setup}

We evaluate \flashOptim with three optimizers: SGD with momentum~\citep{polyak1964some}, AdamW~\citep{loshchilov2017decoupled}, and Lion~\citep{chen2023lion}. We refer to these variants as \flashSGD, \flashAdam, and \flashLion. We test these across several large-scale deep learning tasks, including image classification, language model pretraining, and supervised finetuning.

To ensure a fair comparison, all experiments use identical hyperparameters between reference optimizers and their \flashOptim counterparts. We re-implement the reference optimizers with similar fused Triton kernels for consistent measurement of memory and throughput, and all reference implementations use mixed precision~\citep{micikevicius2018mixed} to keep activations in 16-bit precision. 
The precisions of master weights~$\theta$,
    error correction term~$\rho$, gradients~$g$, momentum~$m$,
    variance~$v$, and activations~$a$ are as follows:

\begin{figure}[!h]
    \centering
    \resizebox{\columnwidth}{!}{%
    \begin{tabular}{lcccccc}
        & $\theta$ & $\rho$ & $g$ & $m$ & $v$ & $a$ \\
        \toprule
        Reference     & FP32 & ---  & FP32 & FP32 & FP32  & BF16 \\
        \flashOptim   & BF16 & INT8 & BF16 & INT8 & UINT8 & BF16 \\
        \bottomrule
    \end{tabular}
    }
\end{figure}


Our experiments demonstrate four main findings.
First, \flashOptim matches reference convergence and accuracy across all tested configurations (\S\ref{sec:convergence}).
Second, it reduces optimizer memory by over 50\% with negligible computational overhead (\S\ref{sec:memory}).
Third, our ULP-based weight splitting achieves near-optimal reconstruction (\S\ref{sec:weight_error}).
Finally, our companding functions significantly reduce quantization error for both momentum and variance states (\S\ref{sec:opt_quant}).


\textbf{Image Classification.} We train a ResNet-50 architecture \citep{he2016resnet} on the ILSVRC2012 (ImageNet-1K) dataset~\citep{deng2009imagenet}. We use the hyperparameters recommended by Nvidia~\citep{nvidia2023resnet}, with additional details provided in Appendix~\ref{sec:appendix_vision}.

\textbf{LLM Pretraining.} We evaluate LLM pretraining using the training recipe outlined in the nanoGPT repository~\citep{karpathy2023nanogpt}.
We use the GPT-2~\citep{radford2019language} architecture and train on a 10B token subset of the FineWeb dataset~\citep{penedo2024fineweb}, following the setup of \citet{jordan2024moddednanogpt}.
We provide hyperparameter details in Appendix~\ref{sec:appendix_llm}.
We evaluate models using a suite of in-context learning (ICL) benchmarks that assess commonsense reasoning and language understanding capabilities. We provide a complete list of benchmarks in Appendix~\ref{sec:appendix_icl}.

\textbf{LLM Finetuning.} We run supervised finetuning on a pretrained Llama-3.1-8B model~\citep{dubey2024llama3} on OpenMathInstruct-2~\citep{toshniwal2024openmathinstruct2}, 
and evaluate on the GSM8k~\citep{cobbe2021gsm8k} benchmark.
We provide further hyperparameter details in Appendix~\ref{sec:appendix_finetuning}.

\textbf{Training and Infrastructure.} We train with distributed data parallelism for the image classification and LLM pretraining tasks, and for LLM finetuning we use FSDP~\citep{zhao2023pytorchfsdp} and activation checkpointing. We train all models using PyTorch 2.8 and CUDA 12.8 on NVIDIA H100 GPUs.
We report the mean and standard deviation for all our results with \nseed random seeds. For loss curve comparisons, we use identical data ordering across methods. System metrics (memory, timing) are measured in steady-state.






\begin{figure}[t]
    \centering

    \begin{subfigure}[t]{\columnwidth}
        \includegraphics[width=\textwidth]{figures/convergence_llm_adamw.pdf}
        \caption{LLM Pretraining (GPT-2 + AdamW)}
        \label{fig:convergence_llm}
    \end{subfigure}
    \begin{subfigure}[t]{\columnwidth}
        \includegraphics[width=\textwidth]{figures/convergence_vision_sgd.pdf}
        \caption{Vision Classification (ResNet-50 + SGD)}
        \label{fig:convergence_vision}
    \end{subfigure}
    \caption{
        \textbf{Training convergence}. 
        Comparison of training loss trajectories between reference optimizers and their \flashOptim variants. Both achieve nearly identical loss curves throughout training, demonstrating that our memory optimizations do not impact model quality.}
    \label{fig:convergence}
\end{figure}


\begin{table}[t]
    \centering
    \caption{
    \textbf{Image Classification \& LLM Finetuning Results}.
    Validation accuracy for ResNet-50 (first two columns), and GSM8k accuracy for the LLM finetuning task (last column). We report standard deviation across \nseed training runs. In all settings \flashOptim matches the reference scores.
    }
    \label{tab:imagenet}
    \small
    \resizebox{\columnwidth}{!}{%

    \begin{tabular}{lccc}
        & \multicolumn{2}{c}{\textbf{ImageNet Top-1 Acc.}} & \textbf{GSM8k Acc.} \\
        \cmidrule(lr){2-3} \cmidrule(lr){4-4}
        & \textbf{SGD} & \textbf{AdamW} & \textbf{AdamW} \\
        \toprule
        Reference    & 77.01 {\scriptsize $\pm$ 0.02} & 75.51 {\scriptsize $\pm$ 0.09} & 75.09 {\scriptsize $\pm$ 0.40} \\
        \flashOptim  & 77.16 {\scriptsize $\pm$ 0.09} & 75.67 {\scriptsize $\pm$ 0.04} & 74.98 {\scriptsize $\pm$ 0.77} \\
        \bottomrule
    \end{tabular}
    }
\end{table}


\begin{table*}
    \centering
    \caption{
    \textbf{LLM Pretraining Results}.
    NanoGPT results with GPT-2 (124M). We report validation loss and accuracy (\%) on in-context learning benchmarks assessing commonsense reasoning and language understanding.
    We report standard deviation across \nseed training runs.
    }
    \label{tab:nanogpt}
    \small
    \resizebox{\textwidth}{!}{%

    \begin{tabular}{lccccccccc}
        \toprule
        \textbf{Optimizer} & \textbf{Val Loss} & \textbf{HellaSwag} & \textbf{ARC-E} & \textbf{CSQA} & \textbf{PIQA} & \textbf{LAMBADA} & \textbf{Winograd} & \textbf{BoolQ} & \textbf{Mean ICL} \\
        \midrule
        AdamW        & 3.263 {\scriptsize $\pm$ 0.001} & 31.9 {\scriptsize $\pm$ 0.2} & 39.6 {\scriptsize $\pm$ 0.7} & 25.9 {\scriptsize $\pm$ 4.1} & 64.3 {\scriptsize $\pm$ 0.0} & 31.0 {\scriptsize $\pm$ 1.2} & 57.3 {\scriptsize $\pm$ 0.6} & 58.1 {\scriptsize $\pm$ 3.6} & \textbf{44.0 {\scriptsize $\pm$ 0.4}} \\
        \flashAdam   & 3.265 {\scriptsize $\pm$ 0.001} & 31.9 {\scriptsize $\pm$ 0.5} & 39.5 {\scriptsize $\pm$ 0.9} & 30.8 {\scriptsize $\pm$ 2.1} & 64.5 {\scriptsize $\pm$ 0.3} & 31.9 {\scriptsize $\pm$ 0.7} & 59.1 {\scriptsize $\pm$ 1.1} & 57.2 {\scriptsize $\pm$ 4.8} & \textbf{45.0 {\scriptsize $\pm$ 1.0}} \\
        \midrule
        Lion         & 3.240 {\scriptsize $\pm$ 0.002} & 32.3 {\scriptsize $\pm$ 0.0} & 40.0 {\scriptsize $\pm$ 0.5} & 23.3 {\scriptsize $\pm$ 1.8} & 63.8 {\scriptsize $\pm$ 1.0} & 31.5 {\scriptsize $\pm$ 0.2} & 58.9 {\scriptsize $\pm$ 2.0} & 58.1 {\scriptsize $\pm$ 2.4} & \textbf{44.0 {\scriptsize $\pm$ 0.5}} \\
        \flashLion   & 3.240 {\scriptsize $\pm$ 0.001} & 32.4 {\scriptsize $\pm$ 0.3} & 40.8 {\scriptsize $\pm$ 0.2} & 25.4 {\scriptsize $\pm$ 2.9} & 64.2 {\scriptsize $\pm$ 0.5} & 31.6 {\scriptsize $\pm$ 0.5} & 59.1 {\scriptsize $\pm$ 2.4} & 59.1 {\scriptsize $\pm$ 2.3} & \textbf{44.7 {\scriptsize $\pm$ 0.5}} \\
        \bottomrule
    \end{tabular}
    }
\end{table*}

\subsection{Convergence and Accuracy} \label{sec:convergence}

We first verify that \flashOptim introduces no measurable degradation by comparing training convergence and validation accuracy.
\autoref{fig:convergence_llm} shows training loss for LLM pretraining with AdamW and \flashAdam.
\flashAdam produces a nearly identical trajectory to the reference AdamW and closely tracks AdamW even after 20,000 parameter updates, indicating that reduced precision does not affect learning dynamics.
\autoref{fig:convergence_vision} shows similar results for image classification: \flashSGD matches reference SGD throughout training.
For LLM finetuning, \autoref{fig:convergence_finetuning} in the appendix shows an analogous result for AdamW. 



\autoref{tab:imagenet} reports final scores for the image classification and LLM finetuning tasks. \flashSGD and \flashAdam match the reference optimizer scores in all three settings.  \autoref{tab:nanogpt} compares final validation loss and in-context learning scores for the LLM pretraining task. 
Models trained with \flashOptim achieve scores within variance of reference optimizers across all metrics.



\subsection{Memory and Speed} \label{sec:memory}

We compare memory requirements and optimizer step time, demonstrating that \flashOptim reduces peak memory~\citep{pytorch_cuda_memory_stats} without practical overhead. We focus on parameter-related memory (weights, optimizer state, gradients) since activation memory is identical across both settings.
To isolate contributions, we ablate: weight splitting (Weight Split) with full-precision optimizer states, and optimizer state quantization with companding (Opt.\ Quant.) with FP32 master weights.

\autoref{tab:finetuning_profiling} breaks down the memory usage for LLM finetuning on a Llama-3.1-8B model. As anticipated, we reduce parameter memory by 50\% from dropping precision from FP32 to BF16, and reduce optimizer memory by 60\% from quantizing the optimizer state tensors. Moreover, when looking at peak memory (including activations), \flashOptim reduces it by 36\% with no practical slowdown in optimizer step times.

Ablating each component confirms that weight splitting halves master weight memory while adding 12\% of extra optimizer state. Optimizer state quantization reduces optimizer state by $\sim$73\% by mapping FP32 tensors to INT8/UINT8; the reduction is slightly less than 75\% due to the overhead of storing FP16 scale factors for each group of 32 elements.
Tables \ref{tab:imagenet_profiling} and \ref{tab:nanogpt_profiling} in the appendix show similar trends for LLM pretraining and image classification.




\begin{table}
    \centering
    \caption{\textbf{Profiling}.
    Runtime measurements for LLM Finetuning with Llama-3.1-8B. We capture master weight memory (Params), optimizer state memory (Optim), peak GPU memory (Peak), and optimizer step times (Step). 
    }
    \label{tab:finetuning_profiling}

    \footnotesize
    \setlength{\tabcolsep}{3pt}
    \begin{tabularx}{\columnwidth}{l *{3}{Yr}c}
        \toprule
        & \multicolumn{2}{c}{\textbf{Params}} & \multicolumn{2}{c}{\textbf{Optim}} & \multicolumn{2}{c}{\textbf{Peak}} & \textbf{Step} \\
        \cmidrule(lr){2-3} \cmidrule(lr){4-5} \cmidrule(lr){6-7} \cmidrule(lr){8-8}
        \textbf{Variant} & GiB & $\Delta$ & GiB & $\Delta$ & GiB & $\Delta$ & ms \\
        \midrule
        Reference    & 29.9 &  & 59.8 &  & 175.2 &  & 12.5 \\
        FlashOptim   & \textbf{15.0} & {\color{black!70}\scriptsize -50\%} & \textbf{23.4} & {\color{black!70}\scriptsize -61\%} & \textbf{112.9} & {\color{black!70}\scriptsize -36\%} & 11.5 \\
        \midrule
        \multicolumn{8}{l}{\emph{Ablations}} \\
        Weight Split & 15.0 & {\color{black!70}\scriptsize -50\%} & 67.3 & {\color{black!70}\scriptsize +12\%} & 156.7 & {\color{black!70}\scriptsize -11\%} & 10.7 \\
        Opt. Quant.  & 29.9 &  & 15.9 & {\color{black!70}\scriptsize -73\%} & 131.4 & {\color{black!70}\scriptsize -25\%} & 10.5 \\
        \bottomrule
    \end{tabularx}
\end{table}

\subsection{Weight Error Correction} \label{sec:weight_error}

We compare our weight splitting scheme to~\citet{zamirai2020revisiting}, who store the rounding error in a floating-point buffer for Kahan summation error correction. Since both approaches are data-independent, we evaluate them exhaustively over all finite FP32 bitstrings, computing relative error after applying compression and decompression in sequence.
We consider four methods: a no error correction baseline, storing error in the same 16-bit format, our ULP-normalized error with 8-bit integers, and ours with 16-bit integers.

\autoref{fig:compression_comparison} plots mean relative error versus exponent for BF16 and FP16.
For BF16, our ULP approach with 16-bit correction achieves near-zero error ($<10^{-9}$).
With 16 bits of correction term, we achieve perfect bitwise reconstruction in $99.92\%$ of the values.
In contrast, storing the error in BF16 (BF16+BF16) produces substantially worse error ($>10^{-6}$), comparable to our 24-bit format.
For FP16, our 32-bit ULP format perfectly reconstructs values in the normal range and dominates FP16+FP16 throughout.
Our 24-bit format produces constant error across the normal range, improving worst-case error from $10^{-4}$ to under $10^{-6}$.

\begin{figure}[t]
    \centering
    \includegraphics[width=\columnwidth]{figures/compression_comparison.pdf}
    \caption{
        \textbf{FP32 Reconstruction Error.}
        Comparison of FP32 reconstruction error for different weight compression schemes across exponent ranges for a target datatype of BF16 (top) and FP16 (bottom).
    Our ULP-based error correction achieves lower relative error particularly for small exponents. 
    Denormal floating point ranges are indicated with vertical dotted lines.
    }
    \label{fig:compression_comparison}
\end{figure}




\subsection{Optimizer State Quantization} \label{sec:opt_quant}

We validate our companding functions by comparing quantization error against standard scaled integer quantization.
Using a fixed full precision training trajectory for consistency, we quantize and dequantize momentum and variance buffers at each step, computing normalized MSE (NMSE) against the original values. 
\autoref{fig:quant_nmse} shows quantile distributions of NMSE for each optimizer and buffer type.
Companding reduces error for momentum buffers and provides substantial improvements for variance buffers, where NMSE drops significantly.

\begin{figure}[t]
    \centering
    \includegraphics[width=\columnwidth]{figures/quant_nmse_compact.pdf}
    \caption{
        \textbf{Optimizer state quantization error.}
        NMSE comparison between standard scaled integer quantization (Linear) and our companding approach for momentum ($m$) and variance ($v$) buffers across different optimizers and datasets.
        Companding reduces quantization error across all optimizer types and tensor types, with particularly large improvements for variance tensors.
    }
    \label{fig:quant_nmse}
\end{figure} 

Beyond reducing quantization error, in some cases companding is essential for training stability.
\autoref{fig:companding_divergence} shows LLM pretraining with and without variance companding: linear quantization causes training to diverge, while companding maintains stable convergence.

\begin{figure}[t]
    \centering
    \includegraphics[width=\columnwidth]{figures/companding_divergence.pdf}
    \caption{
        \textbf{Companding prevents training divergence.}
        GPT-2 training with AdamW and quantized optimizer states: linear quantization (no companding) causes rapid divergence, while our companding approach maintains stable training dynamics.
    }
    \label{fig:companding_divergence}
\end{figure}



\section{Limitations} \label{sec:limitations}



\flashOptim is designed to minimize parameter-associated memory consumption, 
so models with large parameter counts and small activations benefit most from 
our optimizations.
Smaller architectures with large activations, such as convolutional networks 
with high-resolution feature maps, are often dominated by activation memory.
In these activation-dominated regimes, a 50\% reduction in parameter memory 
translates to modest total memory savings and techniques like activation 
checkpointing are more effective for such workloads.

Another limitation is that some tasks and architectures may be more sensitive 
to quantization than those in our benchmarks. The effectiveness of a 
quantization pipeline depends on the data distribution, and there is no guarantee that our method (or any quantization approach) will preserve model quality in all cases. Consequently, our implementation allows selectively disabling compression or excluding specific layers as needed.

Finally, while we find that 24-bit master weights are sufficient in our 
experiments, not even 32 bits are guaranteed to suffice in all cases. Successive (normal) floating-point values differ by a factor of roughly $2^{-\text{num\_mantissa\_bits}}$, and if the ratio of weight update to weight magnitude falls below this threshold, the update will be discarded.



\section{Conclusion} \label{sec:conclusion}


We introduced \flashOptim, a method to reduce the memory footprint of neural 
network training while preserving optimizer semantics and model quality.
\flashOptim provides drop-in replacements for common optimizers and requires 
no additional tuning.

Our approach combines two key techniques. First, we reduce master weights from 32 to 24 bits  via improved floating point error correction. Second, we use companding functions to enable 8-bit optimizer state quantization. Together with 16-bit gradients, these reduce per-parameter memory by over 50\% for AdamW.

We validated \flashOptim on image and language benchmarks using SGD, AdamW, and Lion.
Across all settings, our method matches reference implementations in both loss and accuracy while providing significant memory savings.
\flashOptim composes with FSDP and activation checkpointing, enabling multiplicative benefits for large-scale training.
By lowering memory requirements, \flashOptim enables practitioners and researchers with limited hardware to train larger models than previously feasible.


\section*{Acknowledgements}

We are grateful to Jonathan Frankle, Xing Chen, and Matei Zaharia for their continued support and guidance. We also thank Jialu Liu and Erich Elsen for insightful discussions.


{\small
\begin{thebibliography}{80}
\providecommand{\natexlab}[1]{#1}
\providecommand{\url}[1]{\texttt{#1}}
\expandafter\ifx\csname urlstyle\endcsname\relax
  \providecommand{\doi}[1]{doi: #1}\else
  \providecommand{\doi}{doi: \begingroup \urlstyle{rm}\Url}\fi

\bibitem[Anil et~al.(2019)Anil, Gupta, Koren, and Singer]{anil2019memory}
Rohan Anil, Vineet Gupta, Tomer Koren, and Yoram Singer.
\newblock Memory-efficient adaptive optimization.
\newblock In \emph{Advances in Neural Information Processing Systems
  (NeurIPS)}, 2019.

\bibitem[Biderman et~al.(2024)Biderman, Portes, Ortiz, Paul, Greengard,
  Jennings, King, Havens, Chiley, Frankle, et~al.]{biderman2024lora}
Dan Biderman, Jacob Portes, Jose Javier~Gonzalez Ortiz, Mansheej Paul, Philip
  Greengard, Connor Jennings, Daniel King, Sam Havens, Vitaliy Chiley, Jonathan
  Frankle, et~al.
\newblock Lora learns less and forgets less.
\newblock \emph{arXiv preprint arXiv:2405.09673}, 2024.

\bibitem[Bisk et~al.(2020)Bisk, Zellers, Le~Bras, Gao, and Choi]{bisk2020piqa}
Yonatan Bisk, Rowan Zellers, Ronan Le~Bras, Jianfeng Gao, and Yejin Choi.
\newblock Piqa: Reasoning about physical commonsense in natural language.
\newblock In \emph{Proceedings of the AAAI Conference on Artificial
  Intelligence}, volume~34, pages 7432--7439, 2020.

\bibitem[Brown et~al.(2020)Brown, Mann, Ryder, Subbiah, Kaplan, Dhariwal,
  Neelakantan, Shyam, Sastry, Askell, et~al.]{brown2020language}
Tom Brown, Benjamin Mann, Nick Ryder, Melanie Subbiah, Jared~D Kaplan, Prafulla
  Dhariwal, Arvind Neelakantan, Pranav Shyam, Girish Sastry, Amanda Askell,
  et~al.
\newblock Language models are few-shot learners.
\newblock \emph{Advances in neural information processing systems},
  33:\penalty0 1877--1901, 2020.

\bibitem[Chen et~al.(2016)Chen, Xu, Zhang, and Guestrin]{chen2016training}
Tianqi Chen, Bing Xu, Chiyuan Zhang, and Carlos Guestrin.
\newblock Training deep nets with sublinear memory cost.
\newblock \emph{arXiv preprint arXiv:1604.06174}, 2016.

\bibitem[Chen et~al.(2023)Chen, Liang, Huang, Real, Wang, Liu, Pham, Dong,
  Luong, Hsieh, Lu, and Le]{chen2023lion}
Xiangning Chen, Chen Liang, Da~Huang, Esteban Real, Kaiyuan Wang, Yao Liu, Hieu
  Pham, Xuanyi Dong, Thang Luong, Cho-Jui Hsieh, Yifeng Lu, and Quoc~V. Le.
\newblock Symbolic discovery of optimization algorithms.
\newblock In \emph{Advances in Neural Information Processing Systems
  (NeurIPS)}, 2023.

\bibitem[Chowdhery et~al.(2023)Chowdhery, Narang, Devlin, Bosma, Mishra,
  Roberts, Barham, Chung, Sutton, Gehrmann, et~al.]{chowdhery2023palm}
Aakanksha Chowdhery, Sharan Narang, Jacob Devlin, Maarten Bosma, Gaurav Mishra,
  Adam Roberts, Paul Barham, Hyung~Won Chung, Charles Sutton, Sebastian
  Gehrmann, et~al.
\newblock {PaLM}: Scaling language modeling with pathways.
\newblock \emph{Journal of Machine Learning Research}, 24\penalty0
  (240):\penalty0 1--113, 2023.

\bibitem[Clark et~al.(2019)Clark, Lee, Chang, Kwiatkowski, Collins, and
  Toutanova]{clark2019boolq}
Christopher Clark, Kenton Lee, Ming-Wei Chang, Tom Kwiatkowski, Michael
  Collins, and Kristina Toutanova.
\newblock Boolq: Exploring the surprising difficulty of natural yes/no
  questions.
\newblock In \emph{Proceedings of the Conference of the North American Chapter
  of the Association for Computational Linguistics (NAACL)}, pages 2924--2936,
  2019.

\bibitem[Clark et~al.(2018)Clark, Cowhey, Etzioni, Khot, Sabharwal, Schoenick,
  and Tafjord]{clark2018arc}
Peter Clark, Isaac Cowhey, Oren Etzioni, Tushar Khot, Ashish Sabharwal, Carissa
  Schoenick, and Oyvind Tafjord.
\newblock Think you have solved question answering? try arc, the ai2 reasoning
  challenge.
\newblock \emph{arXiv preprint arXiv:1803.05457}, 2018.

\bibitem[Cobbe et~al.(2021)Cobbe, Kosaraju, Bavarian, Chen, Jun, Kaiser,
  Plappert, Tworek, Hilton, Nakano, Hesse, and Schulman]{cobbe2021gsm8k}
Karl Cobbe, Vineet Kosaraju, Mohammad Bavarian, Mark Chen, Heewoo Jun, Lukasz
  Kaiser, Matthias Plappert, Jerry Tworek, Jacob Hilton, Reiichiro Nakano,
  Christopher Hesse, and John Schulman.
\newblock Training verifiers to solve math word problems.
\newblock \emph{arXiv preprint arXiv:2110.14168}, 2021.

\bibitem[Dehghani et~al.(2023)Dehghani, Djolonga, Mustafa, Padlewski, Heek,
  Gilmer, Steiner, Caron, Geirhos, Alabdulmohsin,
  et~al.]{dehghani2023scalingvit}
Mostafa Dehghani, Josip Djolonga, Basil Mustafa, Piotr Padlewski, Jonathan
  Heek, Justin Gilmer, Andreas~Peter Steiner, Mathilde Caron, Robert Geirhos,
  Ibrahim Alabdulmohsin, et~al.
\newblock Scaling vision transformers to 22 billion parameters.
\newblock In \emph{International Conference on Machine Learning}, pages
  7480--7512. PMLR, 2023.

\bibitem[Deng et~al.(2009)Deng, Dong, Socher, Li, Li, and
  Fei-Fei]{deng2009imagenet}
Jia Deng, Wei Dong, Richard Socher, Li-Jia Li, Kai Li, and Li~Fei-Fei.
\newblock Imagenet: A large-scale hierarchical image database.
\newblock In \emph{2009 IEEE conference on computer vision and pattern
  recognition}, pages 248--255. Ieee, 2009.

\bibitem[Dettmers et~al.(2022)Dettmers, Lewis, Shleifer, and
  Zettlemoyer]{dettmers2022optimizers}
Tim Dettmers, Mike Lewis, Sam Shleifer, and Luke Zettlemoyer.
\newblock 8-bit optimizers via block-wise quantization.
\newblock In \emph{International Conference on Learning Representations
  (ICLR)}, 2022.

\bibitem[Dettmers et~al.(2023)Dettmers, Pagnoni, Holtzman, and
  Zettlemoyer]{dettmers2023qlora}
Tim Dettmers, Artidoro Pagnoni, Ari Holtzman, and Luke Zettlemoyer.
\newblock {QLoRA}: Efficient finetuning of quantized llms.
\newblock \emph{Advances in neural information processing systems},
  36:\penalty0 10088--10115, 2023.

\bibitem[Dubey et~al.(2024)Dubey, Jauhri, Pandey, Kadian, Al-Dahle, Letman,
  Mathur, Schelten, Yang, Fan, et~al.]{dubey2024llama3}
Abhimanyu Dubey, Abhinav Jauhri, Abhinav Pandey, Abhishek Kadian, Ahmad
  Al-Dahle, Aiesha Letman, Akhil Mathur, Alan Schelten, Amy Yang, Angela Fan,
  et~al.
\newblock The llama 3 herd of models.
\newblock \emph{arXiv preprint arXiv:2407.21783}, 2024.

\bibitem[Feng et~al.(2022)Feng, Constable, and Mao]{fsdp2}
Wei Feng, Will Constable, and Yifan Mao.
\newblock Getting started with fully sharded data parallel ({FSDP2}).
\newblock PyTorch Tutorials, 2022.
\newblock URL
  \url{https://docs.pytorch.org/tutorials/intermediate/FSDP_tutorial.html}.
\newblock Last updated: Sep 02, 2025. Accessed: Jan 28, 2026.

\bibitem[Fishman et~al.(2025)Fishman, Chmiel, Banner, and
  Soudry]{fishman2024scaling}
Maxim Fishman, Brian Chmiel, Ron Banner, and Daniel Soudry.
\newblock Scaling {FP8} training to trillion-token {LLM}s.
\newblock In \emph{International Conference on Learning Representations
  (ICLR)}, 2025.
\newblock Spotlight.

\bibitem[Ginsburg et~al.(2019)Ginsburg, Castonguay, Hrinchuk, Kuchaiev,
  Lavrukhin, Leary, Li, Nguyen, Zhang, and Cohen]{ginsburg2019stochastic}
Boris Ginsburg, Patrice Castonguay, Oleksii Hrinchuk, Oleksii Kuchaiev, Vitaly
  Lavrukhin, Ryan Leary, Jason Li, Huyen Nguyen, Yang Zhang, and Jonathan~M
  Cohen.
\newblock Stochastic gradient methods with layer-wise adaptive moments for
  training of deep networks.
\newblock \emph{arXiv preprint arXiv:1905.11286}, 2019.

\bibitem[Goldberg(1991)]{goldberg1991every}
David Goldberg.
\newblock What every computer scientist should know about floating-point
  arithmetic.
\newblock \emph{ACM computing surveys (CSUR)}, 23\penalty0 (1):\penalty0 5--48,
  1991.

\bibitem[{Google}(2019)]{google_cloud_bfloat16_2019}
{Google}.
\newblock Bfloat16: The secret to high performance on cloud tpus.
\newblock Google Cloud Blog, August 2019.

\bibitem[Goyal et~al.(2017)Goyal, Doll{\'a}r, Girshick, Noordhuis, Wesolowski,
  Kyrola, Tulloch, Jia, and He]{goyal2017accurate}
Priya Goyal, Piotr Doll{\'a}r, Ross Girshick, Pieter Noordhuis, Lukasz
  Wesolowski, Aapo Kyrola, Andrew Tulloch, Yangqing Jia, and Kaiming He.
\newblock Accurate, large minibatch sgd: Training imagenet in 1 hour.
\newblock \emph{arXiv preprint arXiv:1706.02677}, 2017.

\bibitem[He et~al.(2016{\natexlab{a}})He, Zhang, Ren, and Sun]{he2016identity}
Kaiming He, Xiangyu Zhang, Shaoqing Ren, and Jian Sun.
\newblock Identity mappings in deep residual networks.
\newblock In \emph{European conference on computer vision}, pages 630--645.
  Springer, 2016{\natexlab{a}}.

\bibitem[He et~al.(2016{\natexlab{b}})He, Zhang, Ren, and Sun]{he2016resnet}
Kaiming He, Xiangyu Zhang, Shaoqing Ren, and Jian Sun.
\newblock Deep residual learning for image recognition.
\newblock In \emph{Proceedings of the IEEE Conference on Computer Vision and
  Pattern Recognition (CVPR)}, pages 770--778, 2016{\natexlab{b}}.

\bibitem[Hoffmann et~al.(2022)Hoffmann, Borgeaud, Mensch, Buchatskaya, Cai,
  Rutherford, Casas, Hendricks, Welbl, Clark, et~al.]{hoffmann2022training}
Jordan Hoffmann, Sebastian Borgeaud, Arthur Mensch, Elena Buchatskaya, Trevor
  Cai, Eliza Rutherford, Diego de~Las Casas, Lisa~Anne Hendricks, Johannes
  Welbl, Aidan Clark, et~al.
\newblock Training compute-optimal large language models.
\newblock \emph{arXiv preprint arXiv:2203.15556}, 2022.

\bibitem[Hu et~al.(2022)Hu, Shen, Wallis, Allen-Zhu, Li, Wang, Wang, and
  Chen]{hu2022lora}
Edward~J. Hu, Yelong Shen, Phillip Wallis, Zeyuan Allen-Zhu, Yuanzhi Li, Shean
  Wang, Lu~Wang, and Weizhu Chen.
\newblock {L}o{RA}: Low-rank adaptation of large language models.
\newblock In \emph{International Conference on Learning Representations
  (ICLR)}, 2022.

\bibitem[Jordan(2024)]{jordan2024moddednanogpt}
Keller Jordan.
\newblock modded-nanogpt: Speedrunning the nanogpt baseline.
\newblock \url{https://github.com/KellerJordan/modded-nanogpt}, 2024.

\bibitem[Jordan et~al.(2024)Jordan, Jin, Boza, You, Cesista, Newhouse, and
  Bernstein]{jordan2024muon}
Keller Jordan, Yuchen Jin, Vlado Boza, Jiacheng You, Franz Cesista, Laker
  Newhouse, and Jeremy Bernstein.
\newblock Muon: An optimizer for hidden layers in neural networks, 2024.
\newblock URL \url{https://kellerjordan.github.io/posts/muon/}.

\bibitem[Kainz et~al.(2004)Kainz, Bogart, and Hess]{kainz2004openexr}
Florian Kainz, Rod Bogart, and Drew Hess.
\newblock The {OpenEXR} image file format.
\newblock In Randima Fernando, editor, \emph{GPU Gems: Programming Techniques,
  Tips and Tricks for Real-Time Graphics}, chapter~26, pages 425--444.
  Addison-Wesley, 2004.

\bibitem[Kalamkar et~al.(2019)Kalamkar, Mudigere, Mellempudi, Das, Banerjee,
  Avancha, Vooturi, Jammalamadaka, Huang, Yuen, et~al.]{kalamkar2019bfloat16}
Dhiraj Kalamkar, Dheevatsa Mudigere, Naveen Mellempudi, Dipankar Das, Kunal
  Banerjee, Sasikanth Avancha, Dharma~Teja Vooturi, Nataraj Jammalamadaka,
  Jianyu Huang, Hector Yuen, et~al.
\newblock A study of bfloat16 for deep learning training.
\newblock \emph{arXiv preprint arXiv:1905.12322}, 2019.

\bibitem[Kaplan et~al.(2020)Kaplan, McCandlish, Henighan, Brown, Chess, Child,
  Gray, Radford, Wu, and Amodei]{kaplan2020scaling}
Jared Kaplan, Sam McCandlish, Tom Henighan, Tom~B Brown, Benjamin Chess, Rewon
  Child, Scott Gray, Alec Radford, Jeffrey Wu, and Dario Amodei.
\newblock Scaling laws for neural language models.
\newblock \emph{arXiv preprint arXiv:2001.08361}, 2020.

\bibitem[Karpathy(2023)]{karpathy2023nanogpt}
Andrej Karpathy.
\newblock nanogpt.
\newblock \url{https://github.com/karpathy/nanoGPT}, 2023.

\bibitem[kjj0(2024)]{kjj0_fineweb10b_gpt2}
kjj0.
\newblock fineweb10b-gpt2.
\newblock Hugging Face Datasets, 2024.
\newblock URL \url{https://huggingface.co/datasets/kjj0/fineweb10B-gpt2}.
\newblock GPT-2 tokenized version of fineweb10B. Accessed: Jan 28, 2026.

\bibitem[Korthikanti et~al.(2023)Korthikanti, Casper, Lym, McAfee, Andersch,
  Shoeybi, and Catanzaro]{korthikanti2023reducing}
Vijay~Anand Korthikanti, Jared Casper, Sangkug Lym, Lawrence McAfee, Michael
  Andersch, Mohammad Shoeybi, and Bryan Catanzaro.
\newblock Reducing activation recomputation in large transformer models.
\newblock \emph{Proceedings of Machine Learning and Systems}, 5:\penalty0
  341--353, 2023.

\bibitem[Krizhevsky et~al.(2012)Krizhevsky, Sutskever, and
  Hinton]{krizhevsky2012imagenet}
Alex Krizhevsky, Ilya Sutskever, and Geoffrey~E Hinton.
\newblock Imagenet classification with deep convolutional neural networks.
\newblock \emph{Advances in neural information processing systems}, 25, 2012.

\bibitem[Levesque et~al.(2012)Levesque, Davis, and
  Morgenstern]{levesque2012winograd}
Hector~J. Levesque, Ernest Davis, and Leora Morgenstern.
\newblock The winograd schema challenge.
\newblock In \emph{Proceedings of the Thirteenth International Conference on
  the Principles of Knowledge Representation and Reasoning (KR)}, pages
  552--561, 2012.

\bibitem[Li et~al.(2023)Li, Chen, and Zhu]{li2023memory}
Bingrui Li, Jianfei Chen, and Jun Zhu.
\newblock Memory efficient optimizers with 4-bit states.
\newblock \emph{Advances in Neural Information Processing Systems},
  36:\penalty0 15136--15171, 2023.

\bibitem[Li and Liang(2021)]{li2021prefix}
Xiang~Lisa Li and Percy Liang.
\newblock Prefix-tuning: Optimizing continuous prompts for generation.
\newblock \emph{arXiv preprint arXiv:2101.00190}, 2021.

\bibitem[Liu et~al.(2025)Liu, Su, Yao, Jiang, Lai, Du, Qin, Xu, Lu, Yan,
  et~al.]{liu2025muon}
Jingyuan Liu, Jianlin Su, Xingcheng Yao, Zhejun Jiang, Guokun Lai, Yulun Du,
  Yidao Qin, Weixin Xu, Enzhe Lu, Junjie Yan, et~al.
\newblock Muon is scalable for llm training.
\newblock \emph{arXiv preprint arXiv:2502.16982}, 2025.

\bibitem[Loshchilov and Hutter(2017)]{loshchilov2017decoupled}
Ilya Loshchilov and Frank Hutter.
\newblock Decoupled weight decay regularization.
\newblock \emph{arXiv preprint arXiv:1711.05101}, 2017.

\bibitem[Lv et~al.(2024{\natexlab{a}})Lv, Yan, Guo, Lv, and Qiu]{lv2024adalomo}
Kai Lv, Hang Yan, Qipeng Guo, Haijun Lv, and Xipeng Qiu.
\newblock Adalomo: Low-memory optimization with adaptive learning rate.
\newblock In \emph{Findings of the Association for Computational Linguistics:
  ACL 2024}, pages 12486--12502, 2024{\natexlab{a}}.

\bibitem[Lv et~al.(2024{\natexlab{b}})Lv, Yang, Liu, Gao, Guo, and
  Qiu]{lv2024lomo}
Kai Lv, Yuqing Yang, Tengxiao Liu, Qinghui Gao, Qipeng Guo, and Xipeng Qiu.
\newblock Full parameter fine-tuning for large language models with limited
  resources.
\newblock In \emph{Proceedings of the 62nd Annual Meeting of the Association
  for Computational Linguistics (ACL)}, pages 8187--8198, 2024{\natexlab{b}}.

\bibitem[Mellempudi et~al.(2019)Mellempudi, Srinivasan, Das, and
  Kaul]{mellempudi2019mixed}
Naveen Mellempudi, Sudarshan Srinivasan, Dipankar Das, and Bharat Kaul.
\newblock Mixed precision training with 8-bit floating point.
\newblock \emph{arXiv preprint arXiv:1905.12334}, 2019.

\bibitem[Micikevicius et~al.(2018)Micikevicius, Narang, Alben, Diamos, Elsen,
  Garcia, Ginsburg, Houston, Kuchaiev, Venkatesh, and
  Wu]{micikevicius2018mixed}
Paulius Micikevicius, Sharan Narang, Jonah Alben, Gregory Diamos, Erich Elsen,
  David Garcia, Boris Ginsburg, Michael Houston, Oleksii Kuchaiev, Ganesh
  Venkatesh, and Hao Wu.
\newblock Mixed precision training.
\newblock In \emph{International Conference on Learning Representations
  (ICLR)}, 2018.

\bibitem[Micikevicius et~al.(2022)Micikevicius, Stosic, Burgess, Cornea, Dubey,
  Grisenthwaite, Ha, Heinecke, Judd, Kamalu, et~al.]{micikevicius2022fp8}
Paulius Micikevicius, Dusan Stosic, Neil Burgess, Marius Cornea, Pradeep Dubey,
  Richard Grisenthwaite, Sangwon Ha, Alexander Heinecke, Patrick Judd, John
  Kamalu, et~al.
\newblock {FP8} formats for deep learning.
\newblock \emph{arXiv preprint arXiv:2209.05433}, 2022.

\bibitem[Mihaylov et~al.(2018)Mihaylov, Clark, Khot, and
  Sabharwal]{mihaylov2018openbookqa}
Todor Mihaylov, Peter Clark, Tushar Khot, and Ashish Sabharwal.
\newblock Can a suit of armor conduct electricity? a new dataset for open book
  question answering.
\newblock In \emph{Proceedings of the Conference on Empirical Methods in
  Natural Language Processing (EMNLP)}, pages 2381--2391, 2018.

\bibitem[Modoranu et~al.(2024)Modoranu, Safaryan, Malinovsky, Kurti{\'c},
  Robert, Richt{\'a}rik, and Alistarh]{modoranu2024microadam}
Ionut-Vlad Modoranu, Mher Safaryan, Grigory Malinovsky, Eldar Kurti{\'c},
  Thomas Robert, Peter Richt{\'a}rik, and Dan Alistarh.
\newblock Micro{A}dam: Accurate adaptive optimization with low space overhead
  and provable convergence.
\newblock \emph{Advances in Neural Information Processing Systems},
  37:\penalty0 1--43, 2024.

\bibitem[{MosaicML}(2022)]{mosaicml2022streaming}
{MosaicML}.
\newblock Streaming: A data loading library for training on large datasets.
\newblock \url{https://github.com/mosaicml/streaming}, 2022.

\bibitem[Narayan et~al.(2025)Narayan, Gupta, Paul, and Blalock]{narayan2025mu}
Saaketh Narayan, Abhay Gupta, Mansheej Paul, and Davis Blalock.
\newblock {$\mu$nit Scaling}: Simple and scalable {FP8} {LLM} training.
\newblock \emph{arXiv preprint arXiv:2502.05967}, 2025.

\bibitem[{NVIDIA}(2023)]{nvidia2023resnet}
{NVIDIA}.
\newblock {ResNet v1.5 for PyTorch}.
\newblock
  \url{https://catalog.ngc.nvidia.com/orgs/nvidia/resources/resnet_50_v1_5_for_pytorch},
  2023.

\bibitem[Paperno et~al.(2016)Paperno, Kruszewski, Lazaridou, Pham, Bernardi,
  Pezzelle, Baroni, Boleda, and Fern{\'a}ndez]{paperno2016lambada}
Denis Paperno, Germ{\'a}n Kruszewski, Angeliki Lazaridou, Quan~Ngoc Pham,
  Raffaella Bernardi, Sandro Pezzelle, Marco Baroni, Gemma Boleda, and Raquel
  Fern{\'a}ndez.
\newblock The lambada dataset: Word prediction requiring a broad discourse
  context.
\newblock In \emph{Proceedings of the 54th Annual Meeting of the Association
  for Computational Linguistics (ACL)}, pages 1525--1534, 2016.

\bibitem[Penedo et~al.(2024)Penedo, Kydl{\'{i}}{\v{c}}ek, Allal, Lozhkov,
  Mitchell, Raffel, Von~Werra, and Wolf]{penedo2024fineweb}
Guilherme Penedo, Hynek Kydl{\'{i}}{\v{c}}ek, Loubna~Ben Allal, Anton Lozhkov,
  Margaret Mitchell, Colin Raffel, Leandro Von~Werra, and Thomas Wolf.
\newblock The fineweb datasets: Decanting the web for the finest text data at
  scale.
\newblock \emph{arXiv preprint arXiv:2406.17557}, 2024.

\bibitem[Peng et~al.(2023)]{peng2023fp8lm}
Houwen Peng et~al.
\newblock Fp8-lm: Training fp8 large language models.
\newblock \emph{arXiv preprint arXiv:2310.18313}, 2023.

\bibitem[Pethick et~al.(2025)Pethick, Xie, Antonakopoulos, Zhu, Silveti-Falls,
  and Cevher]{pethick2025training}
Thomas Pethick, Wanyun Xie, Kimon Antonakopoulos, Zhenyu Zhu, Antonio
  Silveti-Falls, and Volkan Cevher.
\newblock Training deep learning models with norm-constrained lmos.
\newblock In \emph{International Conference on Machine Learning (ICML)}, 2025.

\bibitem[Polyak(1964)]{polyak1964some}
Boris~T Polyak.
\newblock Some methods of speeding up the convergence of iteration methods.
\newblock \emph{Ussr computational mathematics and mathematical physics},
  4\penalty0 (5):\penalty0 1--17, 1964.

\bibitem[{PyTorch Contributors}(2025)]{pytorch_cuda_memory_stats}
{PyTorch Contributors}.
\newblock torch.cuda.memory.memory\_stats.
\newblock PyTorch Documentation, 2025.
\newblock URL
  \url{https://docs.pytorch.org/docs/main/generated/torch.cuda.memory.memory_stats.html}.
\newblock Accessed: Jan 28, 2026.

\bibitem[Radford et~al.(2019)Radford, Wu, Child, Luan, Amodei, Sutskever,
  et~al.]{radford2019language}
Alec Radford, Jeffrey Wu, Rewon Child, David Luan, Dario Amodei, Ilya
  Sutskever, et~al.
\newblock Language models are unsupervised multitask learners.
\newblock \emph{OpenAI blog}, 1\penalty0 (8):\penalty0 9, 2019.

\bibitem[Rajbhandari et~al.(2020)Rajbhandari, Rasley, Ruwase, and
  He]{rajbhandari2020zero}
Samyam Rajbhandari, Jeff Rasley, Olatunji Ruwase, and Yuxiong He.
\newblock Zero: Memory optimizations toward training trillion parameter models.
\newblock In \emph{International Conference for High Performance Computing,
  Networking, Storage and Analysis (SC)}, 2020.

\bibitem[Rajbhandari et~al.(2021)Rajbhandari, Ruwase, Rasley, Smith, and
  He]{rajbhandari2021zeroinfinity}
Samyam Rajbhandari, Olatunji Ruwase, Jeff Rasley, Shaden Smith, and Yuxiong He.
\newblock Zero-infinity: Breaking the gpu memory wall for extreme scale deep
  learning.
\newblock In \emph{International Conference for High Performance Computing,
  Networking, Storage and Analysis (SC)}, 2021.

\bibitem[Ren et~al.(2021)Ren, Rajbhandari, Aminabadi, Ruwase, Yang, Zhang, Li,
  and He]{ren2021zero}
Jie Ren, Samyam Rajbhandari, Reza~Yazdani Aminabadi, Olatunji Ruwase, Shuangyan
  Yang, Minjia Zhang, Dong Li, and Yuxiong He.
\newblock $\{$Zero-offload$\}$: Democratizing $\{$billion-scale$\}$ model
  training.
\newblock In \emph{2021 USENIX Annual Technical Conference (USENIX ATC 21)},
  pages 551--564, 2021.

\bibitem[Rosenfeld et~al.(2019)Rosenfeld, Rosenfeld, Belinkov, and
  Shavit]{rosenfeld2019constructive}
Jonathan~S Rosenfeld, Amir Rosenfeld, Yonatan Belinkov, and Nir Shavit.
\newblock A constructive prediction of the generalization error across scales.
\newblock \emph{arXiv preprint arXiv:1909.12673}, 2019.

\bibitem[Shazeer and Stern(2018)]{shazeer2018adafactor}
Noam Shazeer and Mitchell Stern.
\newblock Adafactor: Adaptive learning rates with sublinear memory cost.
\newblock In \emph{International Conference on Machine Learning (ICML)}, 2018.

\bibitem[Simonyan and Zisserman(2014)]{simonyan2014very}
Karen Simonyan and Andrew Zisserman.
\newblock Very deep convolutional networks for large-scale image recognition.
\newblock \emph{arXiv preprint arXiv:1409.1556}, 2014.

\bibitem[Szegedy et~al.(2016)Szegedy, Vanhoucke, Ioffe, Shlens, and
  Wojna]{szegedy2016rethinking}
Christian Szegedy, Vincent Vanhoucke, Sergey Ioffe, Jon Shlens, and Zbigniew
  Wojna.
\newblock Rethinking the inception architecture for computer vision.
\newblock In \emph{Proceedings of the IEEE conference on computer vision and
  pattern recognition}, pages 2818--2826, 2016.

\bibitem[Talmor et~al.(2019)Talmor, Herzig, Lourie, and
  Berant]{talmor2019commonsenseqa}
Alon Talmor, Jonathan Herzig, Nicholas Lourie, and Jonathan Berant.
\newblock Commonsenseqa: A question answering challenge targeting commonsense
  knowledge.
\newblock In \emph{Proceedings of the Conference of the North American Chapter
  of the Association for Computational Linguistics (NAACL)}, pages 4149--4158,
  2019.

\bibitem[Tan and Le(2019)]{tan2019efficientnet}
Mingxing Tan and Quoc Le.
\newblock Efficientnet: Rethinking model scaling for convolutional neural
  networks.
\newblock In \emph{International conference on machine learning}, pages
  6105--6114. PMLR, 2019.

\bibitem[Tang et~al.(2021)Tang, Gan, Awan, Rajbhandari, Li, Lian, Liu, Zhang,
  and He]{tang20211bit}
Hanlin Tang, Shaoduo Gan, Ammar~Ahmad Awan, Samyam Rajbhandari, Conglong Li,
  Xiangru Lian, Ji~Liu, Ce~Zhang, and Yuxiong He.
\newblock 1-bit {A}dam: Communication efficient large-scale training with
  {A}dam's convergence speed.
\newblock In \emph{International Conference on Machine Learning (ICML)}, 2021.

\bibitem[Tillet et~al.(2019)Tillet, Kung, and Cox]{tillet2019triton}
Philippe Tillet, H.~T. Kung, and David Cox.
\newblock Triton: An intermediate language and compiler for tiled neural
  network computations.
\newblock In \emph{Proceedings of the 3rd ACM SIGPLAN International Workshop on
  Machine Learning and Programming Languages (MAPL)}, pages 10--19, 2019.

\bibitem[Toshniwal et~al.(2024)Toshniwal, Moshkov, Narenthiran, Gitman, Jia,
  and Gitman]{toshniwal2024openmathinstruct2}
Shubham Toshniwal, Ivan Moshkov, Sean Narenthiran, Daria Gitman, Fei Jia, and
  Igor Gitman.
\newblock {OpenMathInstruct-2}: Accelerating ai for math with massive
  open-source instruction data.
\newblock In \emph{Advances in Neural Information Processing Systems
  (NeurIPS)}, 2024.

\bibitem[Vogels et~al.(2019)Vogels, Karimireddy, and Jaggi]{vogels2019powersgd}
Thijs Vogels, Sai~Praneeth Karimireddy, and Martin Jaggi.
\newblock {PowerSGD}: Practical low-rank gradient compression for distributed
  optimization.
\newblock In \emph{Advances in Neural Information Processing Systems
  (NeurIPS)}, 2019.

\bibitem[Wang et~al.(2018)Wang, Choi, Brand, Chen, and
  Gopalakrishnan]{wang2018training}
Naigang Wang, Jungwook Choi, Daniel Brand, Chia-Yu Chen, and Kailash
  Gopalakrishnan.
\newblock Training deep neural networks with 8-bit floating point numbers.
\newblock \emph{Advances in neural information processing systems}, 31, 2018.

\bibitem[Warner(2024)]{warner2024optimi}
Benjamin Warner.
\newblock optimi: Fast, modern, and low precision pytorch optimizers.
\newblock \url{https://github.com/warner-benjamin/optimi}, 2024.

\bibitem[Xi et~al.(2025)Xi, Cai, Zhu, Lu, Keutzer, Chen, and Han]{xi2025coat}
Haocheng Xi, Han Cai, Ligeng Zhu, Yao Lu, Kurt Keutzer, Jianfei Chen, and Song
  Han.
\newblock {COAT}: Compressing optimizer states and activation for
  memory-efficient {FP8} training.
\newblock In \emph{International Conference on Learning Representations
  (ICLR)}, 2025.

\bibitem[Zamirai et~al.(2020)Zamirai, Zhang, Aberger, and
  De~Sa]{zamirai2020revisiting}
Pedram Zamirai, Jian Zhang, Christopher Aberger, and Christopher De~Sa.
\newblock Revisiting bfloat16 training.
\newblock \emph{arXiv preprint arXiv:2010.06192}, 2020.

\bibitem[Zellers et~al.(2019)Zellers, Holtzman, Bisk, Farhadi, and
  Choi]{zellers2019hellaswag}
Rowan Zellers, Ari Holtzman, Yonatan Bisk, Ali Farhadi, and Yejin Choi.
\newblock Hellaswag: Can a machine really finish your sentence?
\newblock In \emph{Proceedings of the 57th Annual Meeting of the Association
  for Computational Linguistics (ACL)}, pages 4791--4800, 2019.

\bibitem[Zhang et~al.(2023)Zhang, Han, Cao, Dai, Miao, Cao, Yang, and
  Xu]{zhang2023adam}
Yijia Zhang, Yibo Han, Shijie Cao, Guohao Dai, Youshan Miao, Ting Cao, Fan
  Yang, and Ningyi Xu.
\newblock Adam accumulation to reduce memory footprints of both activations and
  gradients for large-scale {DNN} training.
\newblock In \emph{European Conference on Artificial Intelligence (ECAI)},
  2023.

\bibitem[Zhang et~al.(2025)Zhang, Chen, Li, Ding, Wu, Ye, Luo, and
  Sun]{zhang2025adammini}
Yushun Zhang, Congliang Chen, Ziniu Li, Tian Ding, Chenwei Wu, Yinyu Ye,
  Zhi-Quan Luo, and Ruoyu Sun.
\newblock Adam-mini: Use fewer learning rates to gain more.
\newblock In \emph{International Conference on Learning Representations
  (ICLR)}, 2025.

\bibitem[Zhao et~al.(2024{\natexlab{a}})Zhao, Zhang, Chen, Wang, Anandkumar,
  and Tian]{zhao2024galore}
Jiawei Zhao, Zhenyu Zhang, Beidi Chen, Zhangyang Wang, Anima Anandkumar, and
  Yuandong Tian.
\newblock {GaLore}: Memory-efficient llm training by gradient low-rank
  projection.
\newblock In \emph{International Conference on Machine Learning (ICML)},
  2024{\natexlab{a}}.

\bibitem[Zhao et~al.(2024{\natexlab{b}})Zhao, Li, Gu, Zheng, K{\"o}lker, Wang,
  and Yuan]{zhao2024adapprox}
Pengxiang Zhao, Ping Li, Yingjie Gu, Yi~Zheng, Stephan~Ludger K{\"o}lker,
  Zhefeng Wang, and Xiaoming Yuan.
\newblock Adapprox: Adaptive approximation in {A}dam optimization via
  randomized low-rank matrices.
\newblock \emph{arXiv preprint arXiv:2403.14958}, 2024{\natexlab{b}}.

\bibitem[Zhao et~al.(2023)Zhao, Gu, Varma, Luo, Huang, Xu, Wright, Shojanazeri,
  Ott, Shleifer, Desmaison, Balioglu, Damania, Nguyen, Chauhan, Hao, Mathews,
  and Li]{zhao2023pytorchfsdp}
Yanli Zhao, Andrew Gu, Rohan Varma, Liang Luo, Chien-Chin Huang, Min Xu, Less
  Wright, Hamid Shojanazeri, Myle Ott, Sam Shleifer, Alban Desmaison, Can
  Balioglu, Pritam Damania, Bernard Nguyen, Geeta Chauhan, Yuchen Hao, Ajit
  Mathews, and Shen Li.
\newblock Pytorch fsdp: Experiences on scaling fully sharded data parallel.
\newblock \emph{Proceedings of the VLDB Endowment (PVLDB)}, 16\penalty0
  (12):\penalty0 3848--3860, 2023.
\newblock \doi{10.14778/3611540.3611569}.

\bibitem[Zhu et~al.(2025)Zhu, Zhang, Cong, Liu, Park, Chandra, Long, Pan, Wang,
  and Lee]{zhu2025apollo}
Hanqing Zhu, Zhenyu Zhang, Wenyan Cong, Xi~Liu, Sem Park, Vikas Chandra,
  Bo~Long, David~Z. Pan, Zhangyang Wang, and Jinwon Lee.
\newblock {APOLLO}: Sgd-like memory, adamw-level performance.
\newblock In \emph{Proceedings of Machine Learning and Systems (MLSys)}, 2025.

\end{thebibliography}
}


\clearpage
\appendix


\section{Method Details}
\label{sec:appendix_method}

This section provides detailed pseudocode for the \flashOptim version of SGD (\autoref{alg:flashsgd}) and Lion (\autoref{alg:flashlion}).

    

\begin{algorithm}[!h]
    \caption{FlashSGD: Memory-Efficient SGD. We \hltext{highlight} the changes from SGD.}
    \label{alg:flashsgd}
    \begin{algorithmic}[1]
    \Require Parameters $\theta_0$, learning rate schedule $\{\eta_t\}_{t=1}^{T}$, momentum $\mu \in [0,1)$, weight decay $\lambda \geq 0$, loss $\mathcal{L}(\theta)$, minibatch sampler $\mathcal{B}(\cdot)$
    \State $\hlchange{m_0^q, m_0^s \gets \mathcal{Q}_m(0)}$
    \State $\hlchange{\thetalp_0, \rho_0 \gets \mathcal{C}(\theta_0)}$
    \State Initialize $t \gets 0$
    \While{not converged}
        \State $t \gets t + 1$
        \State $B_t \sim \mathcal{B}$
        \State $g_t \gets \nabla_\theta \mathcal{L}(B_t; \hlchange{\thetalp_{t-1}})$
        \State $\hlchange{m_{t-1} \gets \mathcal{Q}_m^{-1}(m^q_{t-1}, m^s_{t-1})}$
        \State $m_t \gets \mu m_{t-1} + g_t$
        \State $\hlchange{m^q_t, m^s_t \gets \mathcal{Q}_m(m_t)}$
        \State $\hlchange{\theta_{t-1} \gets \mathcal{C}^{-1}(\thetalp_{t-1}, \rho_{t-1})}$
        \State $\theta_t \gets \theta_{t-1} - \eta_t (m_t + \lambda \theta_{t-1})$
        \State $\hlchange{\thetalp_t, \rho_t \gets \mathcal{C}(\theta_t)}$
    \EndWhile
    \end{algorithmic}
    \end{algorithm}

\begin{algorithm}[!h]
    \caption{FlashLion: Memory-Efficient Lion. We \hltext{highlight} the changes from Lion.}
    \label{alg:flashlion}
    \begin{algorithmic}[1]
    \Require Parameters $\theta_0$, learning rate schedule $\{\eta_t\}_{t=1}^{T}$, $\beta_1, \beta_2 \in [0,1)$, weight decay $\lambda \geq 0$, loss $\mathcal{L}(\theta)$, minibatch sampler $\mathcal{B}(\cdot)$
    \State $\hlchange{m_0^q, m_0^s \gets \mathcal{Q}_m(0)}$
    \State $\hlchange{\thetalp_0, \rho_0 \gets \mathcal{C}(\theta_0)}$
    \State Initialize $t \gets 0$
    \While{not converged}
        \State $t \gets t + 1$
        \State $B_t \sim \mathcal{B}$
        \State $g_t \gets \nabla_\theta \mathcal{L}(B_t; \hlchange{\thetalp_{t-1}})$
        \State $\hlchange{m_{t-1} \gets \mathcal{Q}_m^{-1}(m^q_{t-1}, m^s_{t-1})}$
        \State $u_t \gets \mathrm{sign}(\beta_1 m_{t-1} + (1 - \beta_1) g_t)$ 
        \State $m_t \gets \beta_2 m_{t-1} + (1 - \beta_2) g_t$
        \State $\hlchange{m^q_t, m^s_t \gets \mathcal{Q}_m(m_t)}$
        \State $\hlchange{\theta_{t-1} \gets \mathcal{C}^{-1}(\thetalp_{t-1}, \rho_{t-1})}$
        \State $\theta_t \gets \theta_{t-1} - \eta_t (u_t + \lambda \theta_{t-1})$
        \State $\hlchange{\thetalp_t, \rho_t \gets \mathcal{C}(\theta_t)}$
    \EndWhile
    \end{algorithmic}
    \end{algorithm}






\section{Experimental Details}

For all our experiments we use the MosaicML Streaming~\citep{mosaicml2022streaming} library to ensure deterministic data loading for distributed training.

\subsection{Image Classification}
\label{sec:appendix_vision}

We train the ResNet-50~\citep{he2016resnet} model using the \texttt{timm} library on the ILSVRC2012 (ImageNet-1K) dataset~\citep{deng2009imagenet}, which contains approximately 1.28 million training images and 50,000 validation images across 1,000 classes. Our implementation of ImageNet follows the standard setup from~\citep{krizhevsky2012imagenet,simonyan2014very}. The image is resized with its shorter side randomly sampled in [256, 480] for scale augmentation~\citep{simonyan2014very}. A 224 × 224 crop is randomly sampled from an image or its horizontal flip, and then normalized. For evaluation, the image is first resized to 256 × 256, followed by a 224 × 224 center crop, and then normalized.
We initialize the network with Kaiming He initialization~\citep{he2016identity} and zero-init residuals~\citep{he2016resnet}.

We train for 90 epochs with a batch size of 1024, using a 5-epoch linear warmup followed by cosine learning rate decay, following the recommended settings from~\citep{goyal2017accurate}. We disable weight-decay for biases and BatchNorm layers. We apply label smoothing~\citep{szegedy2016rethinking} with coefficient 0.1. Both reference (FP32 master weights) and \flashOptim use BF16 activations. \autoref{tab:vision_hparams} summarizes the hyperparameters we use for training.

\begin{table}[ht]
\centering
\caption{Optimizer hyperparameters for ImageNet/ResNet-50.}
\label{tab:vision_hparams}
\small
\begin{tabular}{@{}lcc@{}}
\toprule
& SGD & AdamW \\
\midrule
Learning Rate & 1.024 & $3 \times 10^{-3}$ \\
Momentum / Betas & 0.9 & (0.9, 0.999) \\
Weight Decay & $3 \times 10^{-5}$ & $3 \times 10^{-4}$ \\
\bottomrule
\end{tabular}
\end{table}

\paragraph{Memory and Speed Profiling.}
\autoref{tab:imagenet_profiling} presents the results of the memory and speed profile for training.

\begin{table}[ht]
    \centering
    \caption{Memory and speed profiling for image classification (ResNet-50).}
    \label{tab:imagenet_profiling}

    \footnotesize
    \setlength{\tabcolsep}{3pt}
    \begin{tabularx}{\columnwidth}{l *{3}{Yr}c}
        \toprule
        & \multicolumn{2}{c}{\textbf{Params}} & \multicolumn{2}{c}{\textbf{Optim}} & \multicolumn{2}{c}{\textbf{Total}} & \textbf{Step} \\
        \cmidrule(lr){2-3} \cmidrule(lr){4-5} \cmidrule(lr){6-7} \cmidrule(lr){8-8}
        \textbf{Variant} & GiB & $\Delta$ & GiB & $\Delta$ & GiB & $\Delta$ & ms \\
        \midrule
        \multicolumn{8}{l}{\textbf{SGD}} \\
        \quad Reference    & 0.10 &  & 0.10 &  & 0.30 &  & 8.4 \\
        \quad FlashOptim   & 0.05 & {\color{black!70}\scriptsize -46\%} & 0.05 & {\color{black!70}\scriptsize -45\%} & 0.17 & {\color{black!70}\scriptsize -45\%} & 9.0 \\
        \quad Weight Split & 0.05 & {\color{black!70}\scriptsize -46\%} & 0.12 & {\color{black!70}\scriptsize +23\%} & 0.23 & {\color{black!70}\scriptsize -22\%} & 8.7 \\
        \quad Opt. Quant.  & 0.10 &  & 0.03 & {\color{black!70}\scriptsize -73\%} & 0.23 & {\color{black!70}\scriptsize -23\%} & 8.7 \\
        \midrule
        \multicolumn{8}{l}{\textbf{AdamW}} \\
        \quad Reference    & 0.10 &  & 0.19 &  & 0.40 &  & 11.9 \\
        \quad FlashOptim   & 0.05 & {\color{black!70}\scriptsize -46\%} & 0.08 & {\color{black!70}\scriptsize -56\%} & 0.20 & {\color{black!70}\scriptsize -50\%} & 12.2 \\
        \quad Weight Split & 0.05 & {\color{black!70}\scriptsize -46\%} & 0.21 & {\color{black!70}\scriptsize +11\%} & 0.33 & {\color{black!70}\scriptsize -17\%} & 12.1 \\
        \quad Opt. Quant.  & 0.10 &  & 0.05 & {\color{black!70}\scriptsize -73\%} & 0.25 & {\color{black!70}\scriptsize -36\%} & 12.6 \\
        \bottomrule
    \end{tabularx}
\end{table}

\begin{figure}[b]
    \centering
    \includegraphics[width=\columnwidth]{figures/convergence_vision_adamw.pdf}
    \caption{
        \textbf{Training convergence for image classification} (ResNet-50 + AdamW).
        Comparison of validation accuracy between reference AdamW and \flashAdam on ImageNet.
    }
    \label{fig:convergence_vision_adamw}
\end{figure}

\autoref{fig:convergence_vision_adamw} shows training loss for ResNet-50 with AdamW and \flashAdam. \flashAdam produces a nearly identical trajectory to the reference AdamW implementation.

\subsection{LLM Pretraining}
\label{sec:appendix_llm}

We train the GPT-2~\citep{radford2019language} (124M) architecture with 12 transformer layers, 12 attention heads, embedding dimension 768, and context length of 1024. We train on the FineWeb10B~\citep{kjj0_fineweb10b_gpt2} dataset, a subset of approximately 10 billion tokens from the FineWeb~\citep{penedo2024fineweb} dataset, tokenized using the GPT-2 BPE tokenizer. We train for 20,000 steps with a batch size of roughly 0.5 million tokens per step. We apply learning rate warmup for the first 700 steps, followed by a cosine decay to 0. The global norm is clipped at 1.0 and a weight decay of 0.1 is used. Weight decay is applied only to 2D parameters (i.e., weight matrices and embeddings), excluding biases and layer normalization parameters. We train in BF16 mixed precision. \autoref{tab:llm_hparams} summarizes the optimizer hyperparameters.


\begin{table}[ht]
\centering
\caption{Optimizer hyperparameters for GPT-2 124M pretraining.}
\label{tab:llm_hparams}
\small
\begin{tabular}{@{}lcc@{}}
\toprule
& AdamW & Lion \\
\midrule
Learning Rate & $6 \times 10^{-4}$ & $2 \times 10^{-4}$ \\
Betas & (0.9, 0.95) & (0.9, 0.95) \\
\bottomrule
\end{tabular}
\end{table}

\paragraph{Memory and Speed Profiling.}
\autoref{tab:nanogpt_profiling} presents the memory and speed profiling results for LLM pretraining.

\begin{table}[ht]
    \centering
    \caption{Memory and speed profiling for LLM pretraining (GPT-2 124M).}
    \label{tab:nanogpt_profiling}

    \footnotesize
    \setlength{\tabcolsep}{3pt}
    \begin{tabularx}{\columnwidth}{l *{3}{Yr}c}
        \toprule
        & \multicolumn{2}{c}{\textbf{Params}} & \multicolumn{2}{c}{\textbf{Optim}} & \multicolumn{2}{c}{\textbf{Total}} & \textbf{Step} \\
        \cmidrule(lr){2-3} \cmidrule(lr){4-5} \cmidrule(lr){6-7} \cmidrule(lr){8-8}
        \textbf{Variant} & GiB & $\Delta$ & GiB & $\Delta$ & GiB & $\Delta$ & ms \\
        \midrule
        \multicolumn{8}{l}{\textbf{AdamW}} \\
        \quad Reference    & 0.46 &  & 0.93 &  & 1.77 &  & 5.7 \\
        \quad FlashOptim   & 0.23 & {\color{black!70}\scriptsize -50\%} & 0.36 & {\color{black!70}\scriptsize -61\%} & 0.74 & {\color{black!70}\scriptsize -58\%} & 5.9 \\
        \quad Weight Split & 0.23 & {\color{black!70}\scriptsize -50\%} & 1.04 & {\color{black!70}\scriptsize +12\%} & 1.43 & {\color{black!70}\scriptsize -20\%} & 5.9 \\
        \quad Opt. Quant.  & 0.46 &  & 0.25 & {\color{black!70}\scriptsize -73\%} & 1.08 & {\color{black!70}\scriptsize -39\%} & 5.8 \\
        \midrule
        \multicolumn{8}{l}{\textbf{Lion}} \\
        \quad Reference    & 0.46 &  & 0.46 &  & 1.30 &  & 4.3 \\
        \quad FlashOptim   & 0.23 & {\color{black!70}\scriptsize -50\%} & 0.24 & {\color{black!70}\scriptsize -48\%} & 0.62 & {\color{black!70}\scriptsize -53\%} & 4.5 \\
        \quad Weight Split & 0.23 & {\color{black!70}\scriptsize -50\%} & 0.58 & {\color{black!70}\scriptsize +25\%} & 0.96 & {\color{black!70}\scriptsize -26\%} & 4.4 \\
        \quad Opt. Quant.  & 0.46 &  & 0.12 & {\color{black!70}\scriptsize -73\%} & 0.96 & {\color{black!70}\scriptsize -26\%} & 4.4 \\
        \bottomrule
    \end{tabularx}
\end{table}

\begin{figure}[b]
    \centering
    \includegraphics[width=\columnwidth]{figures/convergence_llm_lion.pdf}
    \caption{
        \textbf{Training convergence for LLM pretraining} (GPT-2 + Lion).
        Comparison of validation loss between reference Lion and \flashLion on FineWeb10B.
    }
    \label{fig:convergence_llm_lion}
\end{figure}

\autoref{fig:convergence_llm_lion} shows training loss for LLM pretraining with Lion and \flashLion. \flashLion produces a nearly identical trajectory to the reference Lion implementation and closely tracks Lion even after 20,000 parameter updates.

\subsection{In-Context Learning Benchmarks}
\label{sec:appendix_icl}

We evaluate our pretrained language models on a suite of eight in-context learning~\citep{brown2020language} benchmarks that assess diverse commonsense reasoning and language understanding capabilities:

\begin{itemize}[noitemsep]
    \item \textbf{HellaSwag}~\citep{zellers2019hellaswag}: A sentence completion benchmark requiring physical and temporal commonsense.
    \item \textbf{ARC-Easy}~\citep{clark2018arc}: The easy subset of the AI2 Reasoning Challenge, containing grade-school science questions.
    \item \textbf{CommonsenseQA}~\citep{talmor2019commonsenseqa}: Multiple-choice questions requiring commonsense knowledge from ConceptNet.
    \item \textbf{PIQA}~\citep{bisk2020piqa}: Physical Interaction Question Answering, testing physical commonsense reasoning.
    \item \textbf{OpenBookQA}~\citep{mihaylov2018openbookqa}: Elementary science questions requiring multi-step reasoning over facts.
    \item \textbf{LAMBADA}~\citep{paperno2016lambada}: Word prediction requiring broad discourse context understanding.
    \item \textbf{Winograd}~\citep{levesque2012winograd}: Pronoun resolution problems requiring commonsense reasoning.
    \item \textbf{BoolQ}~\citep{clark2019boolq}: Naturally occurring yes/no reading comprehension questions.
\end{itemize}

All benchmarks are evaluated in a zero-shot setting.

\subsection{LLM Finetuning}
\label{sec:appendix_finetuning}

We fine-tune Llama-3.1-8B~\citep{dubey2024llama3} on OpenMathInstruct-2~\citep{toshniwal2024openmathinstruct2}. For evaluation, we use GSM8k~\citep{cobbe2021gsm8k}, a benchmark of 1,319 grade school math word problems that require multi-step arithmetic reasoning.

\paragraph{Training and Evaluation.}
We use FSDP2~\citep{fsdp2} with full parameter sharding and enabled training with activation checkpointing~\citep{chen2016training} applied to every transformer layer. We use the AdamW optimizer, with $\beta_1=0.9$ and $\beta_2=0.95$. The global norm is clipped at 1.0 and a weight decay of 0.1 is used. Weight decay is applied only to weight matrices, excluding biases, embeddings, and layer normalization parameters. We train for 5000 steps with an effective batch size of approximately 5.2 million tokens per step. We apply learning rate warmup for the first 1000 steps, followed by a cosine decay to 0.

We evaluate on the GSM8k test set using temperature $T=0.2$ decoding.
Following standard practice, we extract the final numerical answer from model generations and compare against ground truth.


\begin{figure}[h]
    \centering
    \includegraphics[width=\columnwidth]{figures/convergence_finetuning_adamw.pdf}
    \caption{
        \textbf{Training convergence for LLM finetuning} (Llama-3.1-8B + AdamW).
        Comparison of training loss between reference AdamW and \flashAdam during supervised finetuning on OpenMathInstruct-2.
    }
    \label{fig:convergence_finetuning}
\end{figure}


\end{document}
